
\documentclass[a4paper,11pt]{article}
\usepackage[utf8]{inputenc}
\usepackage[T1]{fontenc}
\usepackage[spanish]{babel}
\usepackage{amsmath, amssymb, amsfonts}
\usepackage{geometry}
\usepackage{hyperref}
\usepackage{xcolor}
\usepackage{tcolorbox}
\usepackage{booktabs}
\usepackage{listings}
\usepackage{bm}
\usepackage{graphicx}
\usepackage{float}
% TikZ removido (figuras eliminadas)

\geometry{margin=2cm}

\newtcolorbox{concept}[1]{
  colback=blue!5!white,
  colframe=blue!75!black,
  title={\textbf{#1}}
}

\newtcolorbox{implementation}[1]{
  colback=green!5!white,
  colframe=green!60!black,
  title={\textbf{Implementaci\'on: #1}}
}

\newtcolorbox{physics}[1]{
  colback=orange!5!white,
  colframe=orange!70!black,
  title={\textbf{Significado F\'isico: #1}}
}

\newtcolorbox{keyresult}[1]{
  colback=red!5!white,
  colframe=red!70!black,
  title={\textbf{Resultado Clave: #1}}
}

\lstset{
  language=Python,
  basicstyle=\ttfamily\small,
  keywordstyle=\color{blue},
  commentstyle=\color{green!60!black},
  stringstyle=\color{red!70!black},
  frame=single,
  breaklines=true,
  numbers=left,
  numberstyle=\tiny\color{gray}
}

% Colores para cajas de texto (tcolorbox)

\title{\textbf{Retrieval of Aerosol Optical Properties Using\\
Radiative Transfer Theory:\\
A Direct--Inverse Problem Approach}\\[0.5cm]
\large Derivaci\'on Matem\'atica Completa, Soluci\'on Num\'erica e Integraci\'on\\[0.3cm]
\normalsize Basado en el m\'etodo DISORT (Stamnes et al., 1988)}
\author{Emmanuel}
\date{\today}

\begin{document}

\maketitle
\tableofcontents
\newpage

% ======================================================================
\part{Fundamento Te\'orico: La Ecuaci\'on de Transferencia Radiativa}
% ======================================================================

% ======================================================================
\section{Introducci\'on y Motivaci\'on}
% ======================================================================

Este documento presenta la derivaci\'on matem\'atica completa del modelo
\texttt{1D\_Atmosphere\_Model.py}, que resuelve la Ecuaci\'on de Transferencia
Radiativa (RTE) mediante el M\'etodo de Ordenadas Discretas (DISORT) y recupera
propiedades \'opticas de aerosoles mediante m\'ultiples t\'ecnicas de inversi\'on.

\textbf{Objetivo del proyecto:} Dado un perfil atmosf\'erico (espesor \'optico $\tau$,
albedo de dispersi\'on simple $\omega_0$, par\'ametro de asimetr\'ia $g$ por capa,
albedo de superficie $A_s$, \'angulo cenital solar $\theta_0$), calcular las
radiancias ascendentes en el tope de la atm\'osfera (TOA) ---~lo que medir\'ia un
sat\'elite mirando hacia abajo~--- y luego, a partir de esas radiancias observadas,
recuperar los par\'ametros atmosf\'ericos originales.

\textbf{El c\'odigo soporta tres niveles de complejidad:}
\begin{itemize}
    \item \textbf{Nivel 1} --- Atm\'osfera homog\'enea de 1 capa: recuperaci\'on
          del Espesor \'Optico de Aerosoles (AOD) mediante LM+Tikhonov y Optimal
          Estimation.
    \item \textbf{Nivel 2} --- Atm\'osfera realista de 20 capas (US Standard 1976)
          con dispersi\'on de Rayleigh, teor\'ia de Mie para aerosoles, y
          recuperaci\'on conjunta de AOD, SSA y $g$ mediante Optimal Estimation
          (Rodgers, 2000), Phillips--Twomey, y regularizaci\'on Total Variation.
    \item \textbf{Nivel 3} --- Aplicaci\'on a datos reales: simulaci\'on con
          par\'ametros AERONET y recuperaci\'on de AOD a partir de flujos de banda
          ancha medidos por SURFRAD (fracci\'on difusa) mediante bisecci\'on.
\end{itemize}

\textbf{Arquitectura del c\'odigo:}
\begin{itemize}
    \item \texttt{AtmosphereProfile} --- Perfil \'optico multicapa.
    \item \texttt{StandardAtmosphere} --- Atm\'osfera US Standard 1976 con 20 capas y Rayleigh.
    \item \texttt{compute\_aerosol\_optics} --- Teor\'ia de Mie para propiedades de aerosol.
    \item \texttt{DisortSolver} --- Solucionador directo con delta-M scaling y estabilizaci\'on.
    \item \texttt{InverseOptimizer} --- Tikhonov + Levenberg--Marquardt.
    \item \texttt{OptimalEstimation} --- M\'etodo de Rodgers (2000) con diagn\'osticos completos.
    \item \texttt{PhillipsTwomey} --- Inversi\'on lineal restringida con L-curve.
    \item \texttt{TotalVariation} --- Regularizaci\'on TV via IRLS.
\end{itemize}


% [Figura 1: Geometr\'ia plano-paralela -- removida por brevedad]


% ======================================================================
\section{La Ecuaci\'on de Transferencia Radiativa}
% ======================================================================

\subsection{Balance de Energ\'ia en un Camino Infinitesimal}

Un haz de luz con intensidad espectral $I$ [W\,m$^{-2}$\,sr$^{-1}$] recorre
una distancia $ds$ a trav\'es de un medio con part\'iculas.  La variaci\'on de
intensidad se debe a:

\begin{enumerate}
    \item \textbf{Extinci\'on} (p\'erdida): absorci\'on m\'as dispersi\'on fuera
          del haz:  $\;-\sigma_{\text{ext}}\,I\,ds$.
    \item \textbf{Fuente} (ganancia): luz dispersada \emph{hacia} el haz desde
          otras direcciones, m\'as emisi\'on: $\;+\sigma_{\text{ext}}\,J\,ds$.
\end{enumerate}

\begin{equation}\label{eq:rte_path}
    \frac{dI}{ds} = \sigma_{\text{ext}}(J - I)
\end{equation}

\subsection{Coordenadas Verticales y la Variable $\mu$}

En una atm\'osfera plano-paralela, las propiedades dependen solo de la altura
$z$.  Un fot\'on viaja con \'angulo cenital $\theta$ respecto a la vertical, y
la proyecci\'on vertical del camino es $dz = ds\cos\theta$.

\begin{concept}{Variable angular $\mu = \cos\theta$}
\begin{itemize}
    \item $\mu = +1$: direcci\'on vertical ascendente (cenit)
    \item $\mu = -1$: direcci\'on vertical descendente (nadir)
    \item $\mu > 0$: radiaci\'on ascendente (\emph{upwelling});\quad $\mu < 0$: descendente (\emph{downwelling})
    \item $\mu_0 = \cos\theta_0$: coseno del \'angulo cenital solar
\end{itemize}
\end{concept}

Sustituyendo $ds = dz/\mu$:
\begin{equation}
    \mu\,\frac{dI}{dz} = \sigma_{\text{ext}}(J - I)
\end{equation}

\subsection{Espesor \'Optico $\tau$}

\begin{physics}{Espesor \'Optico}
$\tau$ es una coordenada adimensional que mide la extinci\'on acumulada desde el
tope de la atm\'osfera (TOA).  Si $\tau = 1$, el medio aten\'ua la luz directa a
$e^{-1}\approx 37\%$ de su valor original (Ley de Beer--Lambert).
\end{physics}

\begin{equation}
    d\tau = -\sigma_{\text{ext}}\,dz
\end{equation}

El signo negativo asegura que $\tau$ crece desde 0 (TOA) hasta $\tau_{\max}$
(superficie).  Tras la sustituci\'on:

\begin{equation}\label{eq:rte_tau}
    \boxed{\mu\,\frac{dI(\tau,\mu,\phi)}{d\tau} = I(\tau,\mu,\phi) - S(\tau,\mu,\phi)}
\end{equation}

Esta es la ecuaci\'on fundamental (Ec.~1 del reporte DISORT, equivalente a C/I.63
de Chandrasekhar y TS/2.28 de Thomas \& Stamnes).  Aqu\'i $S$ es la funci\'on
fuente total que incluye tanto la dispersi\'on como la emisi\'on.


% ======================================================================
\section{La Funci\'on Fuente y Separaci\'on del Haz Directo}
% ======================================================================

\subsection{Funci\'on Fuente General}

La funci\'on fuente $S$ en la Ec.~\eqref{eq:rte_tau} contiene la dispersi\'on
m\'ultiple desde todas las direcciones m\'as las fuentes internas (Ec.~2a del
reporte DISORT):

\begin{equation}\label{eq:source_general}
    S(\tau,\mu,\phi) = \frac{\omega(\tau)}{4\pi}\int_0^{2\pi}\int_{-1}^{1}
    P(\tau;\cos\Theta)\,I_{\text{tot}}(\tau,\mu',\phi')\,d\mu'\,d\phi'
    + Q(\tau,\mu,\phi)
\end{equation}

donde $P$ es la funci\'on de fase de dispersi\'on, $\omega$ es el albedo de
dispersi\'on simple, y $Q$ es la fuente interna.  El \'angulo de dispersi\'on
$\Theta$ satisface (Ec.~2c del reporte):

\begin{equation}\label{eq:cos_theta}
    \cos\Theta = \mu\mu' + \sqrt{(1-\mu^2)(1-\mu'^2)}\cos(\phi - \phi')
\end{equation}

\subsection{Separaci\'on Difuso--Directo}

El haz solar directo tiene intensidad angular tipo delta de Dirac (Ec.~3b del
reporte DISORT):
\begin{equation}
    I_{\text{direct}} = I_0\,e^{-\tau/\mu_0}\,\delta(\mu+\mu_0)\,\delta(\phi-\phi_0)
\end{equation}

La transformaci\'on difuso-directa (TS/6.2) separa:
\begin{equation}
    I_{\text{tot}} = I_{\text{direct}} + I_{\text{diffuse}}
\end{equation}

Al sustituir en la integral de dispersi\'on y usar la propiedad de la delta de
Dirac ($\int f(x)\delta(x-a)\,dx = f(a)$), se genera un t\'ermino de
\textbf{pseudo-fuente del haz} (Ec.~3d del reporte):

\begin{equation}\label{eq:q_beam}
    Q^{(\text{beam})}(\tau,\mu,\phi) = \frac{\omega(\tau)}{4\pi}\,I_0\,
    P(\tau;\mu,\phi;-\mu_0,\phi_0)\,e^{-\tau/\mu_0}
\end{equation}

\begin{physics}{Por qu\'e se separa el haz directo}
La delta de Dirac no se puede representar con cuadratura num\'erica (ning\'un
punto discreto coincide exactamente con $\mu_0$).  Al evaluarla
anal\'iticamente, se evita un error sistem\'atico.  La parte difusa restante
es suave en $\mu'$ y se aproxima bien con cuadratura.  Esta separaci\'on la
introduce Chandrasekhar (1950) y es est\'andar en DISORT.
\end{physics}

En general, la fuente total es (Ec.~3e del reporte):
\begin{equation}
    Q = Q^{(\text{thermal})} + Q^{(\text{beam})}
\end{equation}

donde $Q^{(\text{thermal})} = (1-\omega)\,B[T(\tau)]$ es la emisi\'on de Planck
(despreciable en el visible/UV).  De aqu\'i en adelante, $I \equiv I_{\text{diffuse}}$
(Ec.~4 del reporte).


% ======================================================================
\section{Funci\'on de Fase y su Expansi\'on en Legendre}
% ======================================================================

\subsection{Separaci\'on de la Dependencia Azimutal}

La funci\'on de fase depende solo del \'angulo de dispersi\'on $\Theta$ (Ec.~2b
del reporte), lo cual permite factorizar la dependencia en $\phi$.  Se expande
$P$ en $2M$ polinomios de Legendre (Ec.~5a del reporte):

\begin{equation}\label{eq:phase_legendre}
    P(\cos\Theta) = \sum_{l=0}^{2M-1}(2l+1)\,g_l\,P_l(\cos\Theta)
\end{equation}

donde $g_l$ son los momentos de Legendre de la funci\'on de fase
(con $g_0 = 1$ por normalizaci\'on).

Usando el teorema de adici\'on de arm\'onicos esf\'ericos (Ec.~5c del reporte):
\begin{equation}
    P_l(\cos\Theta) = P_l^0(\mu)P_l^0(\mu') +
    2\sum_{m=1}^{l}\frac{(l-m)!}{(l+m)!}P_l^m(\mu)P_l^m(\mu')\cos m(\phi-\phi')
\end{equation}

la intensidad y la funci\'on fuente se descomponen en series de Fourier en
$\phi$ (Ec.~6 del reporte para la intensidad; Ec.~8a para la fuente):
\begin{equation}\label{eq:fourier_decomp}
    I(\tau,\mu,\phi) = \sum_{m=0}^{2M-1}I^m(\tau,\mu)\cos m(\phi-\phi_0)
\end{equation}

Cada componente azimutal $m$ satisface una ecuaci\'on independiente (Ec.~7 del reporte):
\begin{equation}\label{eq:rte_azimuthal}
    \mu\,\frac{dI^m}{d\tau} = I^m - \frac{\omega}{2}\sum_{l=m}^{2M-1}
    (2l+1)g_l\,\Lambda_l^m(\mu)\int_{-1}^{1}\Lambda_l^m(\mu')\,I^m(\tau,\mu')\,d\mu'
    - X_0^m(\mu)\,e^{-\tau/\mu_0}
\end{equation}

donde $\Lambda_l^m$ son los polinomios de Legendre asociados normalizados
y $X_0^m$ es el t\'ermino de fuente del haz para la componente $m$.

\begin{concept}{Simplificaci\'on en nuestro c\'odigo}
Nuestro c\'odigo asume \textbf{simetr\'ia azimutal} ($m = 0$ solamente), lo que
significa que la intensidad no depende de $\phi$.  Esto es v\'alido cuando:
\begin{itemize}
    \item La funci\'on de fase es de Henyey--Greenstein (sim\'etrica en $\phi$).
    \item La superficie es Lambertiana (reflexi\'on isotr\'opica).
    \item Solo nos interesan flujos integrados azimutalmente.
\end{itemize}
Con $m=0$, los $\Lambda_l^0(\mu) = P_l(\mu)$ son los polinomios de Legendre
ordinarios, y la RTE se reduce a la Ec.~\eqref{eq:rte_full} m\'as adelante.
\end{concept}

\subsection{Aproximaci\'on de Henyey--Greenstein}

La funci\'on de fase exacta requiere la teor\'ia de Mie (Secci\'on~\ref{sec:mie}).
La parametrizaci\'on de Henyey--Greenstein (HG) la reemplaza con un solo
par\'ametro:

\begin{equation}
    P_{\text{HG}}(\cos\Theta) = \frac{1-g^2}{(1+g^2-2g\cos\Theta)^{3/2}}
\end{equation}

\begin{physics}{Par\'ametro de asimetr\'ia $g$}
\begin{equation}
    g = \langle\cos\Theta\rangle = \frac{1}{2}\int_{-1}^{1}\cos\Theta\,P(\cos\Theta)\,d(\cos\Theta)
\end{equation}
\begin{itemize}
    \item $g=0$: dispersi\'on isotr\'opica (Rayleigh, mol\'eculas de aire)
    \item $g\approx 0.6$--$0.7$: forward scattering dominante (aerosoles t\'ipicos)
    \item $g\to 1$: dispersi\'on casi completamente frontal (gotas grandes de nube)
\end{itemize}
\end{physics}

\subsection{Coeficientes de Legendre: $\chi_l = (2l+1)\,g^l$}

Los coeficientes de la expansi\'on en Legendre de la funci\'on de fase
se obtienen de la \textbf{funci\'on generatriz} de los polinomios de Legendre:

\begin{equation}
    \frac{1}{\sqrt{1 - 2\mu t + t^2}} = \sum_{l=0}^{\infty} P_l(\mu)\,t^l
    \qquad |t| < 1
\end{equation}

Derivando ambos lados y multiplicando por $(1-t^2)$:
\begin{equation}
    \frac{1-t^2}{(1 - 2\mu t + t^2)^{3/2}} = \sum_{l=0}^{\infty}(2l+1)\,P_l(\mu)\,t^l
\end{equation}

Con $t = g$, el lado izquierdo es \emph{exactamente} $P_{\text{HG}}(\mu)$:
\begin{equation}
    P_{\text{HG}}(\mu) = \sum_{l=0}^{\infty}(2l+1)\,g^l\,P_l(\mu)
\end{equation}

\begin{keyresult}{Coeficientes de Henyey--Greenstein}
\begin{equation}
    \boxed{\chi_l = (2l+1)\,g^l}
\end{equation}
Esto no es una aproximaci\'on: es una igualdad exacta para la serie infinita.
En la pr\'actica se trunca en $L = N$ t\'erminos (el n\'umero de streams).
\end{keyresult}

\subsection{Recurrencia de los Polinomios de Legendre}

Los $P_l(\mu)$ se eval\'uan mediante la recurrencia est\'andar:
\begin{equation}
    (l+1)\,P_{l+1}(x) = (2l+1)\,x\,P_l(x) - l\,P_{l-1}(x)
\end{equation}
con $P_0(x) = 1$ y $P_1(x) = x$.  Esta recurrencia es num\'ericamente estable
hasta $l \sim 100$ y tiene complejidad $O(l)$ por punto.


% ======================================================================
\section{RTE Completa en Forma Integral}
% ======================================================================

Combinando la dispersi\'on difusa con la fuente del haz directo, y asumiendo
simetr\'ia azimutal ($m=0$), la RTE completa es:

\begin{equation}\label{eq:rte_full}
    \boxed{\mu\,\frac{dI(\tau,\mu)}{d\tau} = I(\tau,\mu)
    - \frac{\omega_0}{2}\int_{-1}^{1}P(\mu,\mu')\,I(\tau,\mu')\,d\mu'
    - \frac{\omega_0 F_0}{4\pi}\,P(\mu,-\mu_0)\,e^{-\tau/\mu_0}}
\end{equation}

donde la funci\'on de fase como funci\'on de dos cosenos direccionales es:
\begin{equation}\label{eq:phase_matrix_def}
    P(\mu,\mu') = \sum_{l=0}^{N-1}\chi_l\,P_l(\mu)\,P_l(\mu')
\end{equation}

Aqu\'i $I$ se refiere solo a la componente difusa; el haz directo ya fue
extra\'ido.

% [Figura: Radiative balance at optical depth τ -- removida por brevedad]

% ======================================================================
\section{Discretizaci\'on: M\'etodo de Ordenadas Discretas}
% ======================================================================

\subsection{Cuadratura Double-Gauss}

Reemplazamos la integral continua en $\mu'$ por una suma ponderada:
\begin{equation}
    \int_{-1}^{1}f(\mu')\,d\mu' \approx \sum_{j=1}^{N}w_j\,f(\mu_j)
\end{equation}

\textbf{Diferencia con DISORT cl\'asico:} Nuestro c\'odigo utiliza la cuadratura
\textbf{Double-Gauss} (Sykes, 1951; Secci\'on~2.3 del reporte DISORT), en la que
se aplica cuadratura de Gauss-Legendre \emph{por separado} a cada hemisferio
$[0,1]$ y $[-1,0]$, en lugar de usar cuadratura Gauss-Legendre est\'andar sobre
el intervalo completo $[-1,1]$.

Esto implica que para $N$ streams ($N/2$ por hemisferio):
\begin{itemize}
    \item Se calculan $N/2$ nodos $\{\mu_j\}$ y pesos $\{w_j\}$ de GL en $[0,1]$.
    \item Se construyen los $N/2$ nodos descendentes como $\mu_{-j} = -\mu_j$
          con $w_{-j} = w_j$.
\end{itemize}

\begin{concept}{Ventaja del Double-Gauss sobre GL est\'andar}
Como se\~nala el reporte DISORT (Secci\'on~2.3): en la cuadratura GL est\'andar
para $[-1,1]$, los nodos se agrupan solo hacia $\mu = \pm 1$.  Con Double-Gauss,
los nodos tambi\'en se agrupan hacia $\mu = 0$ (\'angulos rasantes), lo que da
resultados superiores en las fronteras donde la intensidad puede ser discontinua
en $\mu = 0$.  Otra ventaja es que los flujos hemisf\'ericos ascendente y
descendente se obtienen directamente sin aproximaciones adicionales.
\end{concept}

\begin{concept}{N\'umero de streams $N$}
Exactamente $N/2$ nodos son positivos (ascendentes) y $N/2$ negativos
(descendentes).  Cada nodo se convierte en una direcci\'on discreta de
propagaci\'on llamada \textbf{stream}.

La precisi\'on de la cuadratura ($2N-1$) debe ser $\geq$ al orden de
truncamiento de la expansi\'on de Legendre ($N$).  Esto se satisface
autom\'aticamente ya que $2N-1 \geq N$ para todo $N \geq 1$.
\end{concept}

% [Figura: Discrete Ordinates & Phase Function -- removida por brevedad]

\subsection{Sistema de EDOs Acopladas}

Evaluando la RTE~\eqref{eq:rte_full} en cada nodo $\mu_i$ y reemplazando la
integral por la cuadratura (Ec.~10a del reporte DISORT; nuestro c\'odigo
trabaja con $m=0$ exclusivamente por la suposici\'on de simetr\'ia azimutal):

\begin{equation}\label{eq:discrete_rte}
    \mu_i\,\frac{dI_i(\tau)}{d\tau} = I_i
    - \frac{\omega_0}{2}\sum_{j=1}^{N}w_j\,P_{ij}\,I_j
    - Q_i\,e^{-\tau/\mu_0}
\end{equation}

donde $I_i(\tau) \equiv I(\tau,\mu_i)$, la \textbf{matriz de fase} es:
\begin{equation}\label{eq:phase_matrix}
    P_{ij} \equiv P(\mu_i,\mu_j) = \sum_{l=0}^{N-1}\chi_l\,P_l(\mu_i)\,P_l(\mu_j)
\end{equation}
y el \textbf{vector de fuente del haz} es:
\begin{equation}\label{eq:beam_source}
    Q_i = \frac{\omega_0 F_0}{4\pi}\,P(\mu_i,-\mu_0)
    = \frac{\omega_0 F_0}{4\pi}\sum_{l=0}^{N-1}\chi_l\,P_l(\mu_i)\,P_l(-\mu_0)
\end{equation}


\subsection{Formulaci\'on Matricial}

Definiendo $\mathbf{I} = [I_1,\ldots,I_N]^T$, el sistema~\eqref{eq:discrete_rte}
se escribe en forma matricial como (forma general multi-stream, Ec.~15 del reporte;
la Ec.~14e muestra la forma equivalente para el caso four-stream con matrices
$\bm{\alpha}$ y $\bm{\beta}$):

\begin{equation}\label{eq:matrix_ode}
    \boxed{\frac{d\mathbf{I}}{d\tau} = \mathbf{A}\,\mathbf{I}
    + \mathbf{S}\,e^{-\tau/\mu_0}}
\end{equation}

Las matrices involucradas son:
\begin{align}
    \mathbf{M} &= \text{diag}(\mu_1,\ldots,\mu_N)
    &\text{(cosenos direccionales)}\label{eq:M_def}\\
    \mathbf{W} &= \text{diag}(w_1,\ldots,w_N)
    &\text{(pesos de cuadratura)}\label{eq:W_def}\\
    \mathbf{P} &= [P_{ij}]_{N\times N}
    &\text{(matriz de fase)}\label{eq:P_def}\\
    \mathbf{A} &= \mathbf{M}^{-1}\!\left(\mathbf{I}_N
    - \tfrac{\omega_0}{2}\,\mathbf{P}\,\mathbf{W}\right)
    &\text{(matriz del sistema)}\label{eq:A_def}\\
    \mathbf{S} &= -\mathbf{M}^{-1}\mathbf{Q}
    &\text{(fuente del haz)}\label{eq:S_def}
\end{align}

\begin{physics}{Estructura de la Matriz $\mathbf{A}$}
Cada fila $i$ de $\mathbf{A}$ divide por $\mu_i$, lo que amplifica las
ecuaciones para direcciones rasantes ($|\mu_i|\ll 1$).  Esto refleja la
f\'isica: un fot\'on con trayectoria casi horizontal recorre un camino
\'optico mucho mayor por unidad de $\Delta\tau$ vertical.  Las filas
correspondientes a $\mu_i > 0$ (ascendentes) y $\mu_i < 0$ (descendentes)
producen eigenvalores de signos opuestos, reflejando la propagaci\'on
en ambos sentidos.
\end{physics}


% ======================================================================
\section{Soluci\'on por Autovalores}
% ======================================================================

\subsection{Soluci\'on Homog\'enea}

La ecuaci\'on~\eqref{eq:matrix_ode} es una EDO lineal no homog\'enea vectorial.
Su soluci\'on general tiene dos partes:

\begin{equation}
    \mathbf{I}(\tau) = \underbrace{\mathbf{I}_h(\tau)}_{\text{homog\'enea}}
    + \underbrace{\mathbf{I}_p(\tau)}_{\text{particular}}
\end{equation}

La ecuaci\'on homog\'enea (sin forzamiento solar) es:
\begin{equation}
    \frac{d\mathbf{I}_h}{d\tau} = \mathbf{A}\,\mathbf{I}_h
\end{equation}

Buscamos soluciones de la forma $\mathbf{I}_h = \mathbf{v}\,e^{-\lambda\tau}$.
Sustituyendo:
\begin{equation}
    -\lambda\,\mathbf{v}\,e^{-\lambda\tau} = \mathbf{A}\,\mathbf{v}\,e^{-\lambda\tau}
    \quad\Longrightarrow\quad
    \mathbf{A}\,\mathbf{v} = -\lambda\,\mathbf{v}
\end{equation}

Esto es la \textbf{ecuaci\'on de autovalores}: $-\lambda$ es un autovalor de
$\mathbf{A}$ y $\mathbf{v}$ su autovector.  En el reporte DISORT, la Ec.~18
presenta este resultado para el caso two-stream; la generalizaci\'on multi-stream
se describe en las Ecs.~23a--f.

Como $\mathbf{A}$ es $N\times N$, tiene $N$ autovalores $\{-\lambda_j\}$ y $N$
autovectores $\{\mathbf{v}_j\}$.  Denotando la matriz de autovectores como
$\mathbf{E} = [\mathbf{v}_1|\cdots|\mathbf{v}_N]$, donde $E_{ij}$ es la componente
$i$ del autovector $j$:

\begin{equation}\label{eq:homogeneous}
    \mathbf{I}_h(\tau) = \sum_{j=1}^{N}C_j\,\mathbf{v}_j\,e^{-\lambda_j\tau}
    \quad\Longleftrightarrow\quad
    I_{h,i}(\tau) = \sum_{j=1}^{N}C_j\,E_{ij}\,e^{-\lambda_j\tau}
\end{equation}

\begin{physics}{Significado f\'isico de los autovalores y autovectores}
\begin{itemize}
    \item Cada par $(\lambda_j, \mathbf{v}_j)$ describe un \textbf{modo natural}
          de transporte radiativo.
    \item $\mathbf{v}_j$ es el \textbf{patr\'on angular}: c\'omo se distribuye
          la intensidad entre los $N$ streams para ese modo.
    \item $\lambda_j$ es la \textbf{tasa de decaimiento}: cu\'an r\'apido cambia
          ese patr\'on con la profundidad \'optica.
    \item $\text{Re}(\lambda_j) > 0$: modo que decae al penetrar
          en la atm\'osfera (componente descendente que se aten\'ua).
    \item $\text{Re}(\lambda_j) < 0$: modo que crece con $\tau$
          (componente ascendente desde la superficie).
    \item $C_j$ es la \textbf{amplitud} de cada modo, determinada por las
          condiciones de frontera.
\end{itemize}
\end{physics}


\subsection{Formulaci\'on Reducida de Autovalores (Ec.~23 del Reporte)}

\textbf{Diferencia con la eigendescomposici\'on directa:} En lugar de resolver
el problema de autovalores completo $N \times N$ de la matriz $\mathbf{A}$,
nuestro c\'odigo implementa la \textbf{formulaci\'on reducida} de Stamnes y
Swanson (1981), que reduce la dimensi\'on a $(N/2) \times (N/2)$, disminuyendo
el costo computacional por un factor $\sim 8$.

La clave es la estructura especial de $\mathbf{A}$: debido a la simetr\'ia
$\mu_{-i} = -\mu_i$ y $w_{-i} = w_i$ (Ec.~14e del reporte), el sistema se
separa en dos subsistemas acoplados para las intensidades ascendentes
$\mathbf{I}^+$ y descendentes $\mathbf{I}^-$:
\begin{equation}
    \frac{d}{d\tau}\begin{bmatrix}\mathbf{I}^+\\\mathbf{I}^-\end{bmatrix}
    = \begin{bmatrix}-\bm{\alpha} & \bm{\beta}\\
    -\bm{\beta} & \bm{\alpha}\end{bmatrix}
    \begin{bmatrix}\mathbf{I}^+\\\mathbf{I}^-\end{bmatrix}
    - \begin{bmatrix}\mathbf{Q}'^+\\\mathbf{Q}'^-\end{bmatrix}
\end{equation}

Buscando soluciones homog\'eneas $\mathbf{I}^\pm = \mathbf{G}^\pm e^{-k\tau}$,
se llega a (Ec.~23c--f del reporte):
\begin{align}
    (\bm{\alpha} - \bm{\beta})(\bm{\alpha} + \bm{\beta})\,\mathbf{v} &= k^2\,\mathbf{v}
    \label{eq:reduced_eigen}
\end{align}

donde $\mathbf{v} = \mathbf{G}^+ + \mathbf{G}^-$.  Este es un problema de
autovalores $(N/2) \times (N/2)$ cuyas soluciones $k_j^2$ son \textbf{reales y
positivas} (Stamnes y Swanson, 1981), por lo que $k_j = \sqrt{k_j^2}$ son
reales.  Los autovectores se recuperan como:
\begin{align}
    \mathbf{G}_j^+ &= \tfrac{1}{2}\left(\mathbf{v}_j + (\bm{\alpha}+\bm{\beta})\,\mathbf{v}_j / k_j\right)\\
    \mathbf{G}_j^- &= \tfrac{1}{2}\left(\mathbf{v}_j - (\bm{\alpha}+\bm{\beta})\,\mathbf{v}_j / k_j\right)
\end{align}

\begin{keyresult}{Ventajas de la formulaci\'on reducida}
\begin{itemize}
    \item \textbf{Autovalores garantizados reales}: No se requiere manejo de
          n\'umeros complejos (todos los arrays son \texttt{float64}).
    \item \textbf{Autovalores en pares $\pm k_j$}: Cada $k_j > 0$ genera
          autom\'aticamente el modo $-k_j$, dando los $N$ modos totales.
    \item \textbf{Reducci\'on $8\times$ en costo}: La eigendescomposici\'on escala
          como $O(n^3)$; pasar de $N$ a $N/2$ da $(N/2)^3 / N^3 = 1/8$.
\end{itemize}
\end{keyresult}


\subsection{Soluci\'on Particular}

Para el t\'ermino forzante $\mathbf{S}\,e^{-\tau/\mu_0}$, proponemos una
soluci\'on particular de la misma forma (Ec.~24a del reporte):
\begin{equation}
    \mathbf{I}_p(\tau) = \mathbf{Z}_p\,e^{-\tau/\mu_0}
\end{equation}

Sustituyendo en~\eqref{eq:matrix_ode}:
\begin{equation}
    -\frac{1}{\mu_0}\,\mathbf{Z}_p\,e^{-\tau/\mu_0}
    = \mathbf{A}\,\mathbf{Z}_p\,e^{-\tau/\mu_0}
    + \mathbf{S}\,e^{-\tau/\mu_0}
\end{equation}

Dividiendo por $e^{-\tau/\mu_0}$ y reorganizando:
\begin{equation}\label{eq:particular}
    \boxed{\left(\mathbf{A} + \frac{1}{\mu_0}\,\mathbf{I}_N\right)\mathbf{Z}_p
    = -\mathbf{S}}
\end{equation}

Este es un sistema lineal $N\times N$ que se resuelve directamente (Ec.~24b
del reporte).

\begin{physics}{Condici\'on de resonancia}
Si $1/\mu_0$ coincide exactamente con un autovalor de $-\mathbf{A}$, la matriz
es singular.  Esto corresponde a una resonancia entre el forzamiento solar y
un modo natural.  Con $N = 16$ streams esto es extremadamente improbable.
\end{physics}

\begin{concept}{Caso especial $\omega_0 = 1$ (dispersi\'on conservativa)}
Cuando $\omega_0 = 1$, el eigenvalor $k = 0$ aparece (Ec.~18b del reporte),
lo que produce singularidades en la formulaci\'on de autovalores.  Nuestro
c\'odigo sigue el enfoque de CDISORT: se fija
$\omega_0 = 1 - 10^{-12}$ para evitar la singularidad, lo que introduce un
error despreciable ($< 10^{-10}$) en la soluci\'on.
\end{concept}


\subsection{Soluci\'on General por Capa}

La soluci\'on completa en una capa con propiedades $(\omega_0, g, \Delta\tau)$ es
(Ec.~27a del reporte):

\begin{equation}\label{eq:general_sol}
    \boxed{I_i(\tau) = \sum_{j=1}^{N}C_j\,E_{ij}\,e^{-\lambda_j\tau}
    + Z_{p,i}\,e^{-\tau/\mu_0}}
\end{equation}

Cada capa tiene sus propios $\mathbf{A}$, eigenvalores, eigenvectores y
$\mathbf{Z}_p$.  Las $N$ constantes $C_j$ por capa se determinan por las
condiciones de frontera.

% [Figura: DISORT Eigenvalue Solution in a Layer -- removida por brevedad]

% ======================================================================
\section{Extensi\'on Multicapa}
% ======================================================================

\subsection{Perfil Atmosf\'erico de $K$ Capas}

La atm\'osfera se divide en $K$ capas (Figura~1 del reporte DISORT), cada una
con propiedades \'opticas independientes $(\Delta\tau_k, \omega_{0,k}, g_k)$.
Las profundidades \'opticas acumuladas en las interfaces son:
\begin{equation}
    \hat{\tau}_0 = 0,\quad
    \hat{\tau}_p = \sum_{k=1}^{p}\Delta\tau_k,\quad
    \hat{\tau}_K = \tau_{\text{total}}
\end{equation}

\subsection{Soluci\'on por Capa con Coordenadas Locales}

En cada capa $k$, se construye su propia matriz $\mathbf{A}_k$ y se
descompone en autovalores:
\begin{align}
    \mathbf{A}_k &= \mathbf{M}^{-1}\!\left(\mathbf{I}_N
    - \tfrac{\omega_{0,k}}{2}\,\mathbf{P}_k\,\mathbf{W}\right)\\
    \mathbf{A}_k\,\mathbf{E}_k &= \mathbf{E}_k\,\bm{\Lambda}_k
    \quad\text{(descomposici\'on espectral)}\\
    \left(\mathbf{A}_k + \mu_0^{-1}\mathbf{I}_N\right)\mathbf{Z}_{p,k} &= -\mathbf{S}_k
\end{align}

La soluci\'on en la capa $k$ se expresa en coordenadas \emph{locales}
$\tau' \in [0, \Delta\tau_k]$ (donde $\tau' = \tau - \hat{\tau}_{k-1}$):

\begin{equation}\label{eq:layer_solution}
    I_i^{(k)}(\tau') = \sum_{j=1}^{N}C_{k,j}\,E^{(k)}_{ij}\,e^{-\lambda^{(k)}_j\tau'}
    + Z_{p,i}^{(k)}\,e^{-(\hat{\tau}_{k-1}+\tau')/\mu_0}
\end{equation}

El sistema tiene $N \times K$ constantes de integraci\'on desconocidas en total.


% ======================================================================
\section{Condiciones de Frontera y Acoplamiento}
% ======================================================================

Las $NK$ constantes se determinan con $NK$ ecuaciones de frontera
(Ecs.~46a--c del reporte):

\subsection{TOA ($\tau=0$): Intensidad Descendente Nula}

No entra radiaci\'on difusa desde el espacio.  Para cada stream descendente
$\mu_i < 0$ ($N/2$ ecuaciones, Ec.~46a del reporte):

\begin{equation}\label{eq:bc_toa}
    I_i^{(1)}(0) = 0 \quad\Longrightarrow\quad
    \sum_{j=1}^{N}C_{1,j}\,E^{(1)}_{ij}\,\underbrace{e^{-\lambda^{(1)}_j\cdot 0}}_{\text{exp estabilizado}}
    = -Z_{p,i}^{(1)}
\end{equation}

\subsection{Continuidad en Interfaces ($\hat{\tau}_p$)}

En cada interfaz entre capas $k$ y $k+1$, la intensidad debe ser continua
para \emph{todos} los $N$ streams ($N$ ecuaciones por interfaz,
$N(K-1)$ en total, Ec.~46b del reporte):

\begin{equation}\label{eq:bc_interface}
    I_i^{(k)}(\Delta\tau_k) = I_i^{(k+1)}(0) \qquad \forall\,i=1,\ldots,N
\end{equation}

Expandiendo ambos lados con la Ec.~\eqref{eq:layer_solution}:
\begin{multline}\label{eq:bc_interface_expanded}
    \sum_j C_{k,j}\,E^{(k)}_{ij}\,e^{-\lambda^{(k)}_j\Delta\tau_k}
    + Z_{p,i}^{(k)}\,e^{-\hat{\tau}_p/\mu_0} \\
    = \sum_j C_{k+1,j}\,E^{(k+1)}_{ij}\,e^{-\lambda^{(k+1)}_j\cdot 0}
    + Z_{p,i}^{(k+1)}\,e^{-\hat{\tau}_p/\mu_0}
\end{multline}

\subsection{BOA ($\tau=\tau_{\text{total}}$): Reflexi\'on Lambertiana}

La superficie refleja con albedo $A_s$.  Para cada stream ascendente
$\mu_i > 0$ ($N/2$ ecuaciones, Ec.~46c del reporte):

\begin{equation}\label{eq:bc_boa}
    I_i^{(K)}(\Delta\tau_K) = \frac{A_s}{\pi}\,F_{\downarrow}(\tau_{\text{total}})
\end{equation}

El flujo descendente total tiene dos componentes:
\begin{equation}\label{eq:flux_down}
    F_{\downarrow} = \underbrace{2\pi\sum_{j:\mu_j<0}w_j|\mu_j|\,I_j^{(K)}(\Delta\tau_K)}_{F_{\text{diff}}\text{ (difuso)}}
    + \underbrace{\mu_0\,F_0\,e^{-\tau_{\text{total}}/\mu_0}}_{F_{\text{dir}}\text{ (directo)}}
\end{equation}

\begin{physics}{Superficie Lambertiana}
Un reflector Lambertiano dispersa la luz incidente uniformemente en todas las
direcciones del hemisferio superior.  El factor $1/\pi$ asegura conservaci\'on
de energ\'ia: $\int_0^{2\pi}\int_0^1 (A_s/\pi)F_{\downarrow}\,\mu\,d\mu\,d\phi
= A_s\,F_{\downarrow}$.
\end{physics}

\subsection{Sistema Global de Frontera}

El conteo de ecuaciones:
\begin{itemize}
    \item TOA: $N/2$ ecuaciones
    \item $(K-1)$ interfaces: $N(K-1)$ ecuaciones
    \item BOA: $N/2$ ecuaciones
    \item \textbf{Total}: $N/2 + N(K-1) + N/2 = NK$ ecuaciones = $NK$ inc\'ognitas $\checkmark$
\end{itemize}

Se ensambla un sistema lineal global:
\begin{equation}\label{eq:global_system}
    \mathbf{B}\,\mathbf{c} = \mathbf{b}
\end{equation}

donde $\mathbf{c} = [C_{1,1},\ldots,C_{1,N},\,C_{2,1},\ldots,C_{K,N}]^T$ es el
vector de $NK$ constantes de integraci\'on.

% [Figura: Multi-Layer Boundary Conditions -- removida por brevedad]

% ======================================================================
\section{Transformaci\'on de Escala (Estabilizaci\'on Num\'erica)}\label{sec:scaling}
% ======================================================================

\subsection{El Problema del Desbordamiento}

La Ec.~\eqref{eq:layer_solution} contiene t\'erminos $e^{-\lambda_j\tau'}$.
Cuando $\text{Re}(\lambda_j) < 0$, el argumento del exponencial se vuelve
\emph{positivo} y la exponencial crece sin l\'imite para $\tau'$ grande.
Esto produce \texttt{NaN} e \texttt{Inf} en la soluci\'on num\'erica.

Como se\~nalan Stamnes y Conklin (1984) y la Secci\'on~2.10 del reporte DISORT,
este mal condicionamiento se elimina completamente mediante una
\textbf{transformaci\'on de escala}.

\subsection{Principio de la Transformaci\'on}

La idea clave (Ec.~48 del reporte) es cambiar las variables de integraci\'on
$C_j$ por variables escaladas $\hat{C}_j$ que absorben las exponenciales
peligrosas:

\begin{itemize}
    \item Para eigenvalores con $\text{Re}(\lambda_j) \geq 0$ (decaen desde el
          tope): se usa el \textbf{tope de la capa} como punto de referencia.
          $\hat{C}_j = C_j$ y el factor exponencial a evaluar es
          $e^{-\lambda_j\tau'}$, que siempre tiene argumento $\leq 0$ para
          $\tau' \geq 0$.

    \item Para eigenvalores con $\text{Re}(\lambda_j) < 0$ (decaen desde el fondo):
          se usa el \textbf{fondo de la capa} como punto de referencia.
          Se define $\hat{C}_j = C_j\,e^{-\lambda_j\Delta\tau}$ y el factor
          exponencial se reescribe como:
          \begin{equation}
              e^{-\lambda_j\tau'} = e^{-\lambda_j(\tau'-\Delta\tau)}
              \cdot e^{-\lambda_j\Delta\tau}
          \end{equation}
          El t\'ermino $e^{-\lambda_j(\tau'-\Delta\tau)}$ tiene argumento
          $= -\lambda_j(\tau'-\Delta\tau) = |\lambda_j|(\Delta\tau - \tau') \geq 0$
          (ya que $\lambda_j < 0$ y $\tau' \leq \Delta\tau$), por lo tanto
          \emph{siempre decae}.
\end{itemize}

\begin{keyresult}{Factor exponencial estabilizado}
El factor exponencial estabilizado en cualquier punto local $\tau'$ dentro de
una capa de espesor $\Delta\tau$ es:
\begin{equation}\label{eq:stable_exp}
    \boxed{f_j(\tau') = \begin{cases}
        e^{-\lambda_j\tau'} & \text{si } \text{Re}(\lambda_j) \geq 0
        \quad\text{(ref.\ al tope)}\\[4pt]
        e^{-\lambda_j(\tau'-\Delta\tau)} & \text{si } \text{Re}(\lambda_j) < 0
        \quad\text{(ref.\ al fondo)}
    \end{cases}}
\end{equation}
En ambos casos, el argumento del exponencial es $\leq 0$, eliminando todo
riesgo de desbordamiento.
\end{keyresult}

\subsection{Verificaci\'on en los L\'imites}

\begin{itemize}
    \item \textbf{Tope de la capa} ($\tau'=0$):
    \begin{itemize}
        \item $\lambda_j \geq 0$: $f_j(0) = e^0 = 1$ \checkmark
        \item $\lambda_j < 0$: $f_j(0) = e^{-\lambda_j(-\Delta\tau)} = e^{|\lambda_j|\Delta\tau \cdot (-1)} = e^{\lambda_j\Delta\tau}$. Como $\lambda_j < 0$, el argumento es negativo \checkmark
    \end{itemize}
    \item \textbf{Fondo de la capa} ($\tau'=\Delta\tau$):
    \begin{itemize}
        \item $\lambda_j \geq 0$: $f_j(\Delta\tau) = e^{-\lambda_j\Delta\tau}$. Argumento $\leq 0$ \checkmark
        \item $\lambda_j < 0$: $f_j(\Delta\tau) = e^0 = 1$ \checkmark
    \end{itemize}
\end{itemize}

En el l\'imite $|\lambda_j|\Delta\tau \to \infty$ (capa \'opticamente gruesa),
la matrix del sistema se reduce a una forma diagonal por bloques, eliminando
completamente el mal condicionamiento (como se demuestra en la Secci\'on~2.10.2
del reporte para el caso two-stream).


% ======================================================================
\section{Delta-M Scaling}\label{sec:deltam}
% ======================================================================

\subsection{El Problema del Pico Frontal}

Para aerosoles con $g > 0.5$, la funci\'on de fase tiene un pico agudo en la
direcci\'on frontal ($\Theta \approx 0$).  Representar este pico con una serie
finita de Legendre truncada en $N$ t\'erminos produce oscilaciones artificiales
(efecto Gibbs) que degradan la precisi\'on.

\subsection{M\'etodo $\delta$-M (Wiscombe, 1977)}

El m\'etodo $\delta$-M (Wiscombe, 1977; Secci\'on~3.2.2 del reporte DISORT,
Ecs.~54a--c) separa la funci\'on de fase en una delta de Dirac frontal
m\'as un residuo suave:

\begin{equation}
    P(\cos\Theta) \approx 2f\,\delta(1-\cos\Theta) + (1-f)\,P'(\cos\Theta)
\end{equation}

donde $f$ es la fracci\'on de energ\'ia en el pico frontal.  Para la
expansi\'on truncada en $2M$ t\'erminos:
\begin{equation}
    f = g_{2M} \approx g^{2M}
\end{equation}

Las propiedades \'opticas escaladas son:
\begin{equation}\label{eq:delta_m_scaling}
    \boxed{
    \begin{aligned}
        \tau' &= \tau\,(1 - \omega_0 f) \\
        \omega_0' &= \frac{\omega_0(1-f)}{1-\omega_0 f} \\
        g' &= \frac{g - f}{1 - f}
    \end{aligned}
    }
\end{equation}

Los coeficientes de Legendre escalados son (Ec.~54b del reporte):
\begin{equation}
    \chi_l' = \frac{\chi_l - f(2l+1)}{1 - f} \qquad l = 0,\ldots,2M-1
\end{equation}

\begin{physics}{Efecto del Delta-M Scaling}
\begin{itemize}
    \item \textbf{Reduce $\tau$}: la extinci\'on efectiva disminuye porque la
          fracci\'on frontal ``no cuenta'' como dispersi\'on real (la luz sigue
          esencialmente en la misma direcci\'on).
    \item \textbf{Reduce $g$}: la funci\'on de fase residual es m\'as sim\'etrica,
          m\'as f\'acil de representar con pocos t\'erminos.
    \item \textbf{Permite usar pocos streams}: con delta-M, $N=16$ streams da
          resultados comparables a $N=64$ sin delta-M para funciones de fase
          con $g > 0.7$.
\end{itemize}
\end{physics}

% [Figura: Delta-M Scaling -- removida por brevedad]

% ======================================================================
\section{Extracci\'on de Radiancias}
% ======================================================================

\subsection{Radiancias TOA Upwelling (Observable del Sat\'elite)}

Una vez resuelto el sistema global~\eqref{eq:global_system} y obtenidas las
constantes $\hat{C}_{k,j}$, evaluamos la intensidad ascendente en $\tau=0$
(tope de la capa 1, $\tau'=0$).  Para cada stream ascendente $\mu_i > 0$:

\begin{equation}\label{eq:I_toa}
    \boxed{I_i^{\text{TOA}} = I_i^{(1)}(\tau'=0)
    = \sum_{j=1}^{N}\hat{C}_{1,j}\,E^{(1)}_{ij}\,f_j(0)
    + Z_{p,i}^{(1)}}
\end{equation}

donde $f_j(0)$ es el factor exponencial estabilizado de la
Ec.~\eqref{eq:stable_exp} evaluado en $\tau'=0$, y el t\'ermino de beam
contribuye con $e^{-0/\mu_0} = 1$.

Esto produce un vector de $N/2$ radiancias en diferentes \'angulos de visi\'on,
que constituye el \textbf{observable} del modelo directo: $\mathbf{y} = F(\mathbf{x})$.

\subsection{Radiancias BOA Downwelling}

De forma an\'aloga, evaluamos en el fondo de la \'ultima capa
($\tau'=\Delta\tau_K$) para streams descendentes $\mu_i < 0$:

\begin{equation}
    I_i^{\text{BOA}} = I_i^{(K)}(\Delta\tau_K)
    = \sum_{j=1}^{N}\hat{C}_{K,j}\,E^{(K)}_{ij}\,f_j(\Delta\tau_K)
    + Z_{p,i}^{(K)}\,e^{-\tau_{\text{total}}/\mu_0}
\end{equation}


% ======================================================================
\section{Intensidades en \'Angulos Arbitrarios (Ecs.~35a,b del Reporte)}
\label{sec:interpolation}
% ======================================================================

Las radiancias obtenidas por el m\'etodo de ordenadas discretas est\'an
disponibles \'unicamente en los $N$ nodos de cuadratura $\{\mu_i\}$.  Para
obtener intensidades en \'angulos \textbf{arbitrarios} $\mu \neq \mu_i$, DISORT
utiliza la \textbf{integraci\'on de la funci\'on fuente} (Stamnes, 1982a;
Secci\'on~2.7--2.8 del reporte).

\subsection{Principio: Integraci\'on Formal de la RTE}

La RTE~\eqref{eq:rte_tau} se puede resolver formalmente para cualquier $\mu > 0$
(ascendente) y $\mu < 0$ (descendente) en un medio de espesor $\tau_L$
(Ec.~28a,b del reporte):
\begin{align}
    I^+(\tau,\mu) &= I^+(\tau_L,\mu)\,e^{-(\tau_L-\tau)/\mu}
    + \int_{\tau}^{\tau_L} S(t,+\mu)\,\frac{e^{-(t-\tau)/\mu}}{\mu}\,dt
    \label{eq:formal_up}\\
    I^-(\tau,\mu) &= I^-(0,\mu)\,e^{-\tau/\mu}
    + \int_{0}^{\tau} S(t,-\mu)\,\frac{e^{-(\tau-t)/\mu}}{\mu}\,dt
    \label{eq:formal_dn}
\end{align}

La funci\'on fuente $S(\tau,\mu)$ se conoce anal\'iticamente a partir de la
soluci\'on discreta (Ec.~30 del reporte):
\begin{equation}\label{eq:source_interp}
    S(\tau,\mu) = \sum_{j=-N}^{N} C_j\,\overline{G}_j(\mu)\,e^{-k_j\tau}
    + \delta_{m0}\left[V_0(\mu) + V_1(\mu)\,\tau\right]
\end{equation}

donde $\overline{G}_j(\mu)$ se calcula interpolando los eigenvectores usando la
cuadratura (Ec.~31a,b del reporte):
\begin{equation}
    \overline{G}_j(\mu) = \sum_{i=-N,\neq 0}^{N} w_i\,D^m(\mu,\mu_i)\,G_j(\mu_i)
\end{equation}

\subsection{Resultado: Ecs.~35a,b del Reporte}

Sustituyendo~\eqref{eq:source_interp} en~\eqref{eq:formal_up}
y~\eqref{eq:formal_dn} e integrando las exponenciales anal\'iticamente capa por
capa (Secci\'on~2.8.2 del reporte), se obtienen las expresiones cerradas para
intensidades en m\'ultiples capas (Ecs.~35a,b).  Para el caso simplificado de
nuestro c\'odigo ($m=0$, sin fuente t\'ermica), la intensidad TOA ascendente
en un \'angulo arbitrario $\mu > 0$ es:
\begin{equation}\label{eq:interp_toa}
    \boxed{I^{\text{TOA}}(\mu) = \sum_{n=1}^{K}\sum_{j=1}^{N}
    C_{jn}\,\frac{\overline{G}_{jn}(\mu)}{1+k_{jn}\mu}\,
    \left(E^+_{jn} - E^+_{jn,\text{next}}\right)
    + \text{t\'erminos de beam}}
\end{equation}

donde los factores $E^+_{jn}$ contienen las exponenciales estabilizadas
que aseguran convergencia num\'erica (Ec.~36 del reporte).

\begin{physics}{Interpolaci\'on vs.\ Radiancias Discretas}
\begin{itemize}
    \item Las radiancias en nodos de cuadratura ($\mu = \mu_i$) son exactas
          al orden de la cuadratura.
    \item Las radiancias interpoladas con Ecs.~35a,b reproducen los valores
          exactos en los nodos de cuadratura (Stamnes, 1982a, hallazgo 1).
    \item Para \'angulos intermedios, las Ecs.~35a,b proporcionan una
          ``correcci\'on inteligente'' que satisface las condiciones de frontera
          y continuidad para \emph{todos} los \'angulos $\mu$ (hallazgo 3).
    \item Nuestro c\'odigo eval\'ua las intensidades en una malla fina de $1°$
          a $89°$ (89 \'angulos) para generar gr\'aficos de campo de radiancia
          continuo, en lugar de los $N/2$ puntos discretos de cuadratura.
\end{itemize}
\end{physics}

\subsection{Uso en la Inversi\'on}

La clase \texttt{InverseOptimizer} acepta un par\'ametro opcional
\texttt{user\_mus} (cosenos de \'angulos de visi\'on definidos por el usuario).
Cuando se proporciona, el modelo directo $F(\mathbf{x})$ utiliza la
interpolaci\'on de Ecs.~35a,b en lugar de devolver las intensidades en los
nodos de cuadratura:
\begin{equation}
    F(\mathbf{x}) = \begin{cases}
        [I(\mu_1), \ldots, I(\mu_{N/2})]^T & \text{sin \texttt{user\_mus}}\\
        [I(\mu_1^u), \ldots, I(\mu_M^u)]^T & \text{con \texttt{user\_mus}}
    \end{cases}
\end{equation}

Esto permite ajustar el modelo a observaciones en \'angulos que no coinciden
con los nodos de cuadratura, como ocurrir\'ia con datos reales de sat\'elite.


% ======================================================================
\part{Modelo Atmosf\'erico Realista}
% ======================================================================

% ======================================================================
\section{Atm\'osfera Est\'andar US-1976}\label{sec:us76}
% ======================================================================

Para el Nivel~2, se construye una atm\'osfera realista de 20 capas basada en la
US Standard Atmosphere 1976.

\subsection{Estructura Vertical}

21 fronteras de capa (20 capas), m\'as densas cerca de la superficie donde
se concentran los aerosoles:
\begin{equation}
    z = [0, 1, 2, 3, 4, 5, 6, 8, 10, 12, 15, 20, 25, 30, 35, 40, 50, 60, 70, 85, 100]\;\text{km}
\end{equation}

\subsection{Perfil de Temperatura}

La temperatura sigue segmentos lineales (\emph{lapse rates}):
\begin{center}
\begin{tabular}{llr}
\toprule
\textbf{Regi\'on} & \textbf{Altitud} & \textbf{$L$ (K/km)} \\ \midrule
Troposfera      & 0--11 km   & $-6.5$ \\
Tropopausa      & 11--20 km  & $0$ \\
Estratosfera    & 20--32 km  & $+1.0$ \\
                & 32--47 km  & $+2.8$ \\
Estratopausa    & 47--51 km  & $0$ \\
Mesosfera       & 51--71 km  & $-2.8$ \\
                & 71--100 km & $-2.0$ \\
\bottomrule
\end{tabular}
\end{center}

\subsection{Presi\'on: Ecuaci\'on Hidrost\'atica}

\begin{equation}
    P(z) = \begin{cases}
        P(z_0)\left(\dfrac{T(z)}{T(z_0)}\right)^{-g_0 M_{\text{air}}/(R\,L)}
        & L \neq 0 \\[8pt]
        P(z_0)\exp\!\left(-\dfrac{g_0 M_{\text{air}}(z-z_0)}{R\,T(z_0)}\right)
        & L = 0
    \end{cases}
\end{equation}

con $g_0 = 9.80665$~m/s$^2$, $M_{\text{air}} = 0.0289644$~kg/mol,
$R = 8.31447$~J/(mol$\cdot$K).

\subsection{Densidad Num\'erica}

\begin{equation}
    n_{\text{air}} = \frac{P}{k_B\,T} \qquad\text{[m}^{-3}\text{]}
\end{equation}

% [Figura: 20-Layer Atmosphere Profile -- removida por brevedad]

% ======================================================================
\section{Dispersi\'on de Rayleigh}\label{sec:rayleigh}
% ======================================================================

\subsection{Secci\'on Eficaz de Rayleigh}

La secci\'on eficaz de dispersi\'on Rayleigh por mol\'ecula
(Bodhaine et al., 1999):
\begin{equation}
    \sigma_{\text{Ray}} = \frac{4.02\times 10^{-28}}{\lambda_{\mu\text{m}}^{4.04}}
    \quad\text{[cm}^2\text{]}
\end{equation}

\subsection{Espesor \'Optico de Rayleigh por Capa}

\begin{equation}
    \tau_{\text{Ray},k} = \sigma_{\text{Ray}} \cdot \bar{n}_{\text{air},k}
    \cdot \Delta z_k
\end{equation}

donde $\bar{n}_{\text{air},k}$ es la densidad num\'erica media en la capa $k$
y $\Delta z_k$ es el espesor geom\'etrico.

\subsection{Funci\'on de Fase de Rayleigh}

\begin{equation}
    P_{\text{Ray}}(\cos\Theta) = \frac{3}{4}(1 + \cos^2\Theta)
\end{equation}

En Legendre: $\chi_0 = 1$, $\chi_2 = 0.5$, todos los dem\'as $= 0$.
Es decir, $g_{\text{Ray}} = 0$ (dispersi\'on sim\'etrica).

\subsection{Distribuci\'on Vertical del Aerosol}\label{sec:aer_vertical}

El AOD total de la columna ($\tau_{\text{aer,total}}$) se distribuye
verticalmente seg\'un un \textbf{perfil exponencial decreciente} con una
altura de escala $H_{\text{aer}}$ (en nuestro c\'odigo, $H_{\text{aer}} = 2.0$~km):
\begin{equation}\label{eq:aer_profile}
    \tau_{\text{aer},k} = \tau_{\text{aer,total}} \cdot
    \frac{e^{-z_{\text{mid},k} / H_{\text{aer}}}}
    {\sum_{j=1}^{K} e^{-z_{\text{mid},j} / H_{\text{aer}}}}
\end{equation}

donde $z_{\text{mid},k}$ es la altitud media de la capa $k$
(promedio de sus fronteras superior e inferior).  Este perfil concentra la
mayor\'ia del aerosol en las primeras capas:
\begin{itemize}
    \item A $z = 0$~km: peso m\'aximo $\propto 1$.
    \item A $z = H_{\text{aer}} = 2$~km: peso $\propto e^{-1} \approx 0.37$.
    \item A $z = 5$~km: peso $\propto e^{-2.5} \approx 0.08$.
    \item Por encima de 10~km: contribuci\'on despreciable ($< 1\%$).
\end{itemize}

\subsection{Propiedades Combinadas Rayleigh + Aerosol}

Para cada capa $k$, las propiedades \'opticas totales se obtienen combinando
las contribuciones de Rayleigh (molecular) y aerosol.  La dispersi\'on de
Rayleigh es puramente conservativa ($\omega_{\text{Ray}} = 1$) e isotr\'opica
($g_{\text{Ray}} = 0$):

\begin{align}
    \tau_{\text{total},k} &= \tau_{\text{aer},k} + \tau_{\text{Ray},k}
    \label{eq:tau_combined}\\
    \omega_{0,k} &= \frac{\tau_{\text{aer},k}\,\omega_{\text{aer}}
    + \tau_{\text{Ray},k}\cdot 1}{\tau_{\text{total},k}}
    \label{eq:ssa_combined}\\
    g_{\text{eff},k} &= \frac{\tau_{\text{aer},k}\,\omega_{\text{aer}}\,g_{\text{aer}}
    + \tau_{\text{Ray},k}\cdot 1 \cdot 0}
    {\tau_{\text{aer},k}\,\omega_{\text{aer}} + \tau_{\text{Ray},k}}
    \label{eq:g_combined}
\end{align}

\begin{physics}{Significado de las f\'ormulas de combinaci\'on}
\begin{itemize}
    \item \textbf{Ec.~\eqref{eq:tau_combined}}: El espesor \'optico total es
          aditivo (extinciones se suman).
    \item \textbf{Ec.~\eqref{eq:ssa_combined}}: El SSA combinado pondera por
          el espesor \'optico de cada componente.  En capas altas donde domina
          Rayleigh ($\tau_{\text{Ray}} \gg \tau_{\text{aer}}$),
          $\omega_{0,k} \to 1$ (dispersi\'on pura).  En capas bajas con mucho
          aerosol, $\omega_{0,k} \to \omega_{\text{aer}}$ (contribuci\'on de
          absorci\'on del aerosol).
    \item \textbf{Ec.~\eqref{eq:g_combined}}: El par\'ametro de asimetr\'ia
          efectivo pondera por el \emph{espesor \'optico de dispersi\'on}
          ($\tau \cdot \omega$) de cada componente.  Solo la dispersi\'on
          del aerosol contribuye asimetr\'ia ($g > 0$); Rayleigh aporta
          dispersi\'on sim\'etrica ($g = 0$).
    \item El SSA combinado se acota: $\omega_{0,k} \in (0, 1-10^{-12}]$ para
          evitar la singularidad de dispersi\'on conservativa.
\end{itemize}
\end{physics}


% ======================================================================
\section{Teor\'ia de Mie para Propiedades del Aerosol}\label{sec:mie}
% ======================================================================

\subsection{Eficiencias de Mie}

Para una esfera de radio $r$ con \'indice de refracci\'on complejo
$m = n_r + in_i$, el par\'ametro de tama\~no es:
\begin{equation}
    x = \frac{2\pi r}{\lambda}
\end{equation}

Las eficiencias de extinci\'on y dispersi\'on se calculan mediante las series
de Mie (Bohren \& Huffman, 1983):
\begin{align}
    Q_{\text{ext}} &= \frac{2}{x^2}\sum_{n=1}^{n_{\max}}(2n+1)\,\text{Re}(a_n + b_n)\\
    Q_{\text{sca}} &= \frac{2}{x^2}\sum_{n=1}^{n_{\max}}(2n+1)(|a_n|^2 + |b_n|^2)
\end{align}

donde $a_n$ y $b_n$ son los coeficientes de Mie (Bohren \& Huffman, 1983,
Cap.~4):
\begin{align}\label{eq:mie_an_bn}
    a_n &= \frac{\bigl[D_n(mx)/m + n/x\bigr]\,\psi_n(x) - \psi_{n-1}(x)}
               {\bigl[D_n(mx)/m + n/x\bigr]\,\xi_n(x)  - \xi_{n-1}(x)}\\[6pt]
    b_n &= \frac{\bigl[D_n(mx)\cdot m + n/x\bigr]\,\psi_n(x) - \psi_{n-1}(x)}
               {\bigl[D_n(mx)\cdot m + n/x\bigr]\,\xi_n(x)  - \xi_{n-1}(x)}
\end{align}

Aqu\'i:
\begin{itemize}
    \item $m = n_r + i\,n_i$ es el \'indice de refracci\'on complejo del
          aerosol relativo al medio circundante.
    \item $D_n(z) = \frac{d}{dz}\ln\psi_n(z) = \frac{\psi_n'(z)}{\psi_n(z)}$
          es la \textbf{derivada logar\'itmica} de la funci\'on de Riccati--Bessel,
          calculada por \textbf{recurrencia descendente} (Wiscombe, 1980):
          \begin{equation}
              D_{n-1}(z) = \frac{n}{z} - \frac{1}{D_n(z) + n/z}
          \end{equation}
          partiendo de $D_{n_{\text{mx}}} = 0$ con
          $n_{\text{mx}} = \max(n_{\max}, |mx|) + 16$.
    \item $\psi_n(x)$ y $\chi_n(x)$ son las funciones de Riccati--Bessel,
          generadas por recurrencia ascendente:
          \begin{equation}
              f_{n+1}(x) = \frac{2n+1}{x}\,f_n(x) - f_{n-1}(x)
          \end{equation}
          con condiciones iniciales $\psi_0 = \sin x$,
          $\psi_1 = \sin x/x - \cos x$,
          $\chi_0 = \cos x$, $\chi_1 = \cos x/x + \sin x$.
    \item $\xi_n(x) = \psi_n(x) - i\,\chi_n(x)$ es la funci\'on de
          Riccati--Bessel de tercer tipo.
    \item $n_{\max}$ se determina con el criterio de Wiscombe (1980):
          $n_{\max} = x + 4\,x^{1/3} + 2$.
\end{itemize}

El par\'ametro de asimetr\'ia $g$ se obtiene de:
\begin{equation}
    g = \frac{4}{Q_{\text{sca}}\,x^2}\sum_{n=2}^{n_{\max}}
    \left[\frac{(n-1)(n+1)}{n}\text{Re}(a_{n-1}a_n^* + b_{n-1}b_n^*)
    + \frac{2n-1}{(n-1)n}\text{Re}(a_{n-1}b_{n-1}^*)\right]
\end{equation}

\subsection{Integraci\'on sobre Distribuci\'on de Tama\~nos}

Los aerosoles atmosf\'ericos siguen una distribuci\'on lognormal:
\begin{equation}
    n(r) = \frac{N_0}{r\,\ln\sigma_g\,\sqrt{2\pi}}\,
    \exp\!\left[-\frac{(\ln r - \ln r_g)^2}{2\ln^2\sigma_g}\right]
\end{equation}

Las propiedades \'opticas \emph{bulk} (promediadas) son:
\begin{align}
    B_{\text{ext}} &= \int_0^{\infty} Q_{\text{ext}}(r)\,\pi r^2\,n(r)\,dr\\
    B_{\text{sca}} &= \int_0^{\infty} Q_{\text{sca}}(r)\,\pi r^2\,n(r)\,dr\\
    \text{SSA} &= B_{\text{sca}} / B_{\text{ext}}\\
    g_{\text{bulk}} &= \frac{\int_0^{\infty} g(r)\,Q_{\text{sca}}(r)\,\pi r^2\,n(r)\,dr}{B_{\text{sca}}}
\end{align}


% ======================================================================
\part{El Problema Inverso}
% ======================================================================

% [Figura: Forward vs Inverse Problem -- removida por brevedad]

% ======================================================================
\section{Planteamiento General del Problema Inverso}
% ======================================================================

\subsection{Definici\'on del Problema}

El \textbf{problema directo} (forward problem) consiste en calcular las
radiancias TOA $\mathbf{y}$ a partir de par\'ametros atmosf\'ericos
$\mathbf{x}$ mediante el solver DISORT:
\begin{equation}
    \mathbf{y} = F(\mathbf{x}) + \bm{\epsilon}
\end{equation}
donde:
\begin{itemize}
    \item $F: \mathbb{R}^n \to \mathbb{R}^m$ es el \textbf{modelo directo}
          (operador no lineal que mapea par\'ametros a observaciones).
    \item $\mathbf{x} \in \mathbb{R}^n$ es el \textbf{vector de estado}
          (par\'ametros a recuperar).
    \item $\mathbf{y} \in \mathbb{R}^m$ es el \textbf{vector de observaci\'on}
          ($m = N/2$ radiancias TOA en los nodos de cuadratura ascendentes).
    \item $\bm{\epsilon} \sim \mathcal{N}(\mathbf{0}, \mathbf{S}_\epsilon)$ es
          el \textbf{error de medici\'on} (ruido gaussiano).
\end{itemize}

El \textbf{problema inverso} busca recuperar $\mathbf{x}$ a partir de
$\mathbf{y}_{\text{obs}}$:
\begin{itemize}
    \item \textbf{Nivel 1}: $\mathbf{x} = [\text{AOD}]^T$ ($n=1$, 1 inc\'ognita)
    \item \textbf{Nivel 2}: $\mathbf{x} = [\text{AOD},\,\omega_0,\,g]^T$
          ($n=3$, 3 inc\'ognitas)
\end{itemize}

\subsection{Mal Condicionamiento (Ill-Posedness)}

Seg\'un Hadamard (1902), un problema est\'a \textbf{bien planteado} si la
soluci\'on existe, es \'unica, y depende continuamente de los datos.
Nuestro problema inverso viola las tres condiciones:
\begin{enumerate}
    \item \textbf{No unicidad}: Diferentes combinaciones de $(\omega_0, g)$
          producen radiancias similares (degeneraci\'on AOD--SSA,
          Secci\'on~\ref{sec:interp_coupling}).
    \item \textbf{Inestabilidad}: Peque\~nas perturbaciones en
          $\mathbf{y}_{\text{obs}}$ causan grandes cambios en $\mathbf{x}$
          (amplificaci\'on del ruido).
    \item \textbf{Subdeterminaci\'on}: $m = N/2$ observaciones angulares deben
          determinar par\'ametros con informaci\'on redundante limitada.
\end{enumerate}

La \textbf{regularizaci\'on} resuelve el mal condicionamiento a\~nadiendo
informaci\'on \emph{a priori} sobre la soluci\'on esperada, estabilizando el
problema a costa de introducir un sesgo (bias) controlado.

\subsection{Cuatro M\'etodos Implementados}

Nuestro c\'odigo implementa cuatro m\'etodos de regularizaci\'on, cada uno
con un enfoque diferente:
\begin{center}
\small
\begin{tabular}{@{}lp{4cm}p{4cm}l@{}}
\toprule
\textbf{M\'etodo} & \textbf{Enfoque} & \textbf{Uso en nuestro c\'odigo}
& \textbf{Nivel} \\ \midrule
LM + Tikhonov & Penalizaci\'on $L^2$ escalar & Recuperaci\'on principal
& 1, 2 \\
Optimal Estimation & Bayesiano (covarianzas) & Recuperaci\'on + diagn\'osticos
& 1, 2 \\
Phillips--Twomey & Inversi\'on lineal + L-curve & Demo de selecci\'on de $\gamma$
& 2 \\
Total Variation & Penalizaci\'on $L^1$ (IRLS) & Disponible, no usado
& --- \\
\bottomrule
\end{tabular}
\end{center}

Los dos primeros (LM+Tikhonov y OE) se usan para la recuperaci\'on real de
par\'ametros.  Phillips--Twomey se usa para demostrar el m\'etodo de la
L-curve.  Total Variation est\'a implementado para uso futuro (perfiles con
discontinuidades).

% ======================================================================
\section{Regularizaci\'on de Tikhonov con Levenberg--Marquardt}
\label{sec:tikhonov}
% ======================================================================

\subsection{Funci\'on de Costo}

El m\'etodo de Tikhonov (Tikhonov, 1963) a\~nade un t\'ermino de
penalizaci\'on que mantiene la soluci\'on cerca de un estado de referencia
$\mathbf{x}_a$ (llamado \textbf{a priori} o \textbf{estado de fondo}):

\begin{equation}\label{eq:cost_tikhonov}
    \boxed{J(\mathbf{x}) = \underbrace{\|\mathbf{y}_{\text{obs}} - F(\mathbf{x})\|^2}_{\text{ajuste a datos}}
    + \underbrace{\gamma\,\|\mathbf{x} - \mathbf{x}_a\|^2}_{\text{regularizaci\'on}}}
\end{equation}

donde cada s\'imbolo se define como:
\begin{itemize}
    \item $\|\mathbf{v}\|^2 = \mathbf{v}^T\mathbf{v} = \sum_i v_i^2$:
          norma euclidiana al cuadrado.
    \item $F(\mathbf{x})$: modelo directo DISORT que calcula las $m$ radiancias
          TOA para el estado $\mathbf{x}$.
    \item $\mathbf{x}_a$: \textbf{estado a priori}, una estimaci\'on inicial
          razonable de la soluci\'on.  En nuestro c\'odigo:
          $\mathbf{x}_a = [0.5]$ (Nivel 1) o $[0.5, 0.85, 0.5]$ (Nivel 2).
    \item $\gamma > 0$: \textbf{par\'ametro de regularizaci\'on} (escalar).
          Controla el balance entre ajustar los datos y mantenerse cerca de
          $\mathbf{x}_a$:
          \begin{itemize}
              \item $\gamma \to 0$: la soluci\'on sigue solo los datos
                    (puede amplificar el ruido).
              \item $\gamma \to \infty$: la soluci\'on $\to \mathbf{x}_a$
                    (ignora los datos).
              \item Nuestro c\'odigo usa $\gamma = 0.01$ (Nivel~1) y
                    $\gamma = 0.05$ (Nivel~2).
          \end{itemize}
\end{itemize}

\subsection{Residuo Aumentado}

Para usar \texttt{scipy.optimize.least\_squares} (que minimiza
$\|\mathbf{r}\|^2$ para un vector residuo $\mathbf{r}$), se reformula
$J(\mathbf{x})$ como un problema de m\'inimos cuadrados \textbf{aumentado}:
\begin{equation}\label{eq:augmented_residual}
    \mathbf{r}(\mathbf{x}) = \begin{bmatrix}
        \mathbf{y}_{\text{obs}} - F(\mathbf{x}) \\[4pt]
        \sqrt{\gamma}\,(\mathbf{x} - \mathbf{x}_a)
    \end{bmatrix}
    \in \mathbb{R}^{m+n}
\end{equation}

El truco es que al calcular $\|\mathbf{r}\|^2$:
\begin{equation}
    \|\mathbf{r}\|^2 = \|\mathbf{y}_{\text{obs}} - F(\mathbf{x})\|^2
    + \gamma\,\|\mathbf{x} - \mathbf{x}_a\|^2 = J(\mathbf{x})
\end{equation}

El vector $\mathbf{r}$ tiene $m+n$ componentes: las primeras $m$ miden el
desajuste con los datos, y las \'ultimas $n$ act\'uan como
``pseudo-observaciones'' que atraen la soluci\'on hacia $\mathbf{x}_a$.

\subsection{Algoritmo de Levenberg--Marquardt (LM)}

El algoritmo LM (Levenberg, 1944; Marquardt, 1963) es un m\'etodo iterativo
para minimizar $\|\mathbf{r}(\mathbf{x})\|^2$.  En cada iteraci\'on $n$,
el \textbf{Jacobiano} del residuo se define como:
\begin{equation}
    \mathbf{J}_n = \frac{\partial\mathbf{r}}{\partial\mathbf{x}}\bigg|_{\mathbf{x}_n}
    \in \mathbb{R}^{(m+n)\times n}
\end{equation}

que nuestro c\'odigo calcula por \textbf{diferencias finitas}:
\begin{equation}
    [\mathbf{J}_n]_{ij} \approx \frac{r_i(\mathbf{x}_n + \delta_j\mathbf{e}_j) - r_i(\mathbf{x}_n)}{\delta_j}
\end{equation}

donde $\mathbf{e}_j$ es el $j$-\'esimo vector unitario y
$\delta_j = \max(|x_j| \cdot 10^{-2},\, 10^{-8})$.

La actualizaci\'on LM resuelve:
\begin{equation}\label{eq:lm_update}
    (\mathbf{J}_n^T\mathbf{J}_n + \lambda_{\text{LM}}\,\mathbf{I})\,\Delta\mathbf{x}
    = -\mathbf{J}_n^T\,\mathbf{r}(\mathbf{x}_n)
\end{equation}

donde $\lambda_{\text{LM}} \geq 0$ es el \textbf{par\'ametro de amortiguamiento}
que controla el tama\~no del paso:
\begin{itemize}
    \item $\lambda_{\text{LM}} \to 0$: paso de \textbf{Gauss--Newton}
          (converge r\'apido cerca del m\'inimo pero puede divergir lejos).
    \item $\lambda_{\text{LM}} \to \infty$: paso de \textbf{descenso de gradiente}
          (robusto pero lento).
\end{itemize}

\texttt{scipy.optimize.least\_squares} ajusta $\lambda_{\text{LM}}$
autom\'aticamente en cada iteraci\'on usando la estrategia Trust Region
Reflective (TRF) cuando se proporcionan cotas para los par\'ametros, o LM
puro cuando no hay cotas.

\begin{physics}{Sesgo de Tikhonov}
El t\'ermino $\gamma\|\mathbf{x}-\mathbf{x}_a\|^2$ introduce un
\textbf{sesgo sistem\'atico}: la soluci\'on se desplaza hacia $\mathbf{x}_a$.
En el Nivel~1 con $\mathbf{x}_a = [0.5]$, los AOD verdaderos extremos
($\tau < 0.1$ o $\tau > 1.0$) se recuperan con error porque el regularizador
los atrae hacia 0.5.  Este efecto se reduce con $\gamma$ peque\~no, pero
no se elimina completamente.  La Optimal Estimation (Secci\'on~\ref{sec:oe})
permite un control m\'as refinado de este sesgo mediante covarianzas.
\end{physics}


% ======================================================================
\section{Optimal Estimation (Rodgers, 2000)}\label{sec:oe}
% ======================================================================

\subsection{Fundamento Bayesiano}

La estimaci\'on \'optima (OE) trata el problema inverso como uno de
\textbf{inferencia bayesiana} (Rodgers, 2000, Cap.~2--5).  Se asume que tanto
el estado verdadero como las mediciones son variables aleatorias con
distribuciones gaussianas conocidas:

\begin{itemize}
    \item \textbf{A priori}: $\mathbf{x} \sim \mathcal{N}(\mathbf{x}_a, \mathbf{S}_a)$,
          donde $\mathbf{x}_a \in \mathbb{R}^n$ es la mejor estimaci\'on del
          estado antes de medir, y $\mathbf{S}_a \in \mathbb{R}^{n\times n}$
          es su \textbf{matriz de covarianza a priori}.
          \begin{itemize}
              \item $[\mathbf{S}_a]_{ii} = \sigma_{a,i}^2$: varianza del
                    par\'ametro $i$ (incertidumbre al cuadrado).
              \item $[\mathbf{S}_a]_{ij}$ ($i\neq j$): covarianza entre
                    par\'ametros (t\'ipicamente cero si son independientes).
              \item En nuestro c\'odigo (Nivel~2):
                    $\mathbf{S}_a = \text{diag}(0.3^2, 0.1^2, 0.2^2)$,
                    lo que significa $\sigma_a(\text{AOD}) = 0.3$,
                    $\sigma_a(\omega_0) = 0.1$, $\sigma_a(g) = 0.2$.
          \end{itemize}
    \item \textbf{Medici\'on}: $\mathbf{y} = F(\mathbf{x}) + \bm{\epsilon}$
          con $\bm{\epsilon} \sim \mathcal{N}(\mathbf{0}, \mathbf{S}_\epsilon)$,
          donde $\mathbf{S}_\epsilon \in \mathbb{R}^{m\times m}$ es la
          \textbf{matriz de covarianza del error de medici\'on}.
          \begin{itemize}
              \item En nuestro c\'odigo:
                    $\mathbf{S}_\epsilon = \sigma_\epsilon^2\,\mathbf{I}_m$
                    con $\sigma_\epsilon = 0.02\,\bar{y}$ (2\% de la
                    radiancia media), asumiendo errores no correlacionados
                    entre \'angulos.
          \end{itemize}
\end{itemize}

Por el \textbf{teorema de Bayes}, la probabilidad posterior del estado dado
las mediciones es:
\begin{equation}
    P(\mathbf{x}|\mathbf{y}) \propto P(\mathbf{y}|\mathbf{x})\,P(\mathbf{x})
\end{equation}

Sustituyendo las distribuciones gaussianas:
\begin{equation}
    P(\mathbf{x}|\mathbf{y}) \propto
    \exp\!\left[-\frac{1}{2}\left(
    \underbrace{(\mathbf{y}-F(\mathbf{x}))^T\mathbf{S}_\epsilon^{-1}(\mathbf{y}-F(\mathbf{x}))}_{\chi^2_{\text{datos}}}
    + \underbrace{(\mathbf{x}-\mathbf{x}_a)^T\mathbf{S}_a^{-1}(\mathbf{x}-\mathbf{x}_a)}_{\chi^2_{\text{a priori}}}
    \right)\right]
\end{equation}

Maximizar la probabilidad posterior equivale a minimizar la \textbf{funci\'on
de costo}:
\begin{equation}\label{eq:cost_oe}
    \boxed{J(\mathbf{x}) = (\mathbf{y}-F(\mathbf{x}))^T\mathbf{S}_\epsilon^{-1}(\mathbf{y}-F(\mathbf{x}))
    + (\mathbf{x}-\mathbf{x}_a)^T\mathbf{S}_a^{-1}(\mathbf{x}-\mathbf{x}_a)}
\end{equation}

\begin{concept}{Relaci\'on con Tikhonov}
Comparando las Ecs.~\eqref{eq:cost_tikhonov} y~\eqref{eq:cost_oe}: Tikhonov
usa un escalar $\gamma$ que penaliza todos los par\'ametros por igual; OE usa
matrices $\mathbf{S}_a^{-1}$ y $\mathbf{S}_\epsilon^{-1}$ que permiten pesar
cada par\'ametro y cada medici\'on de forma independiente.  Si
$\mathbf{S}_a = \gamma^{-1}\mathbf{I}$ y $\mathbf{S}_\epsilon = \mathbf{I}$,
OE se reduce exactamente a Tikhonov.
\end{concept}

\subsection{Iteraci\'on de Gauss--Newton}

Como $F(\mathbf{x})$ es no lineal, se minimiza $J$ iterativamente.  En cada
iteraci\'on $n$, se linealiza el modelo directo alrededor del estado actual
$\mathbf{x}_n$:
\begin{equation}
    F(\mathbf{x}) \approx F(\mathbf{x}_n) + \mathbf{K}_n(\mathbf{x}-\mathbf{x}_n)
\end{equation}

donde $\mathbf{K}_n \in \mathbb{R}^{m \times n}$ es el \textbf{Jacobiano}
(o \textbf{weighting function matrix}):
\begin{equation}
    [\mathbf{K}_n]_{ij} = \frac{\partial F_i}{\partial x_j}\bigg|_{\mathbf{x}_n}
\end{equation}

Nuestro c\'odigo calcula $\mathbf{K}$ por \textbf{diferencias finitas
hacia adelante}: se perturba cada par\'ametro $x_j$ una cantidad
$\delta_j = \max(|x_j| \cdot 0.01, 10^{-8})$ y se mide el cambio en las
radiancias.

Sustituyendo la linealizaci\'on en~\eqref{eq:cost_oe} y derivando con
respecto a $\mathbf{x}$, la actualizaci\'on es (Rodgers, 2000, Ec.~5.8):

\begin{equation}\label{eq:oe_update}
    \boxed{\mathbf{x}_{n+1} = \mathbf{x}_a + \mathbf{G}_n\left[
    \mathbf{y} - F(\mathbf{x}_n) + \mathbf{K}_n(\mathbf{x}_n - \mathbf{x}_a)
    \right]}
\end{equation}

donde la \textbf{gain matrix} (matriz de ganancia, Rodgers Ec.~5.9) es:
\begin{equation}\label{eq:gain}
    \mathbf{G}_n = \underbrace{\left(\mathbf{K}_n^T\mathbf{S}_\epsilon^{-1}\mathbf{K}_n
    + \mathbf{S}_a^{-1}\right)^{-1}}_{\hat{\mathbf{S}}_n}\mathbf{K}_n^T\mathbf{S}_\epsilon^{-1}
\end{equation}

La expresi\'on entre par\'entesis invertida es la \textbf{covarianza posterior}
$\hat{\mathbf{S}}_n$, que ya contiene toda la informaci\'on sobre la
incertidumbre de la soluci\'on.

Opcionalmente, nuestro c\'odigo a\~nade amortiguamiento LM al t\'ermino
de la inversa: $\hat{\mathbf{S}}_n^{-1} + \gamma_{\text{LM}}\,
\text{diag}(\hat{\mathbf{S}}_n^{-1})$, lo que estabiliza las primeras
iteraciones lejos del m\'inimo (usado en Nivel~2b con $\gamma_{\text{LM}} = 0.05$).

\subsection{Diagn\'osticos de la Recuperaci\'on}

Una ventaja fundamental de OE sobre Tikhonov es que proporciona
\textbf{diagn\'osticos cuantitativos} de la calidad de la inversi\'on.

\subsubsection{Covarianza Posterior $\hat{\mathbf{S}}$}

\begin{equation}
    \hat{\mathbf{S}} = \left(\mathbf{K}^T\mathbf{S}_\epsilon^{-1}\mathbf{K}
    + \mathbf{S}_a^{-1}\right)^{-1}
\end{equation}

$\hat{\mathbf{S}}$ es la covarianza de la distribuci\'on posterior gaussiana.
$\sqrt{[\hat{\mathbf{S}}]_{ii}}$ da la \textbf{incertidumbre posterior} del
par\'ametro $i$: si los datos son informativos, $\sqrt{[\hat{\mathbf{S}}]_{ii}}
< \sigma_{a,i}$ (la incertidumbre se reduce respecto al a priori).

\subsubsection{Averaging Kernel $\mathbf{A}$}

\begin{equation}
    \mathbf{A} = \mathbf{G}\,\mathbf{K} = \frac{\partial\hat{\mathbf{x}}}{\partial\mathbf{x}_{\text{true}}}
    \in \mathbb{R}^{n \times n}
\end{equation}

El \textbf{averaging kernel} describe la sensibilidad de la soluci\'on
recuperada al estado verdadero.  Cada fila $i$ de $\mathbf{A}$ muestra c\'omo
el par\'ametro recuperado $\hat{x}_i$ depende de todos los par\'ametros
verdaderos:
\begin{itemize}
    \item $\mathbf{A} = \mathbf{I}_n$: recuperaci\'on perfecta (cada par\'ametro
          se determina independientemente).
    \item $\mathbf{A} \approx \mathbf{0}$: la soluci\'on est\'a dominada por
          el a priori (los datos no aportan informaci\'on).
\end{itemize}

\subsubsection{Degrees of Freedom for Signal (DFS)}

\begin{equation}
    \text{DFS} = \text{tr}(\mathbf{A})
\end{equation}

DFS es el n\'umero efectivo de piezas de informaci\'on independiente
extra\'idas de las mediciones.  Para $n$ par\'ametros:
$\text{DFS} \leq n$.  En nuestro Nivel~2 ($n=3$), DFS $\approx 2.5$--$2.8$
indica que no toda la informaci\'on de los 3 par\'ametros es independiente
(existe acoplamiento AOD--SSA--$g$).

\subsubsection{Criterio de Convergencia}

La convergencia se eval\'ua con la \textbf{distancia de Mahalanobis}
del paso (Rodgers, 2000, Ec.~5.29):
\begin{equation}
    d^2 = (\mathbf{x}_{n+1}-\mathbf{x}_n)^T\hat{\mathbf{S}}_n^{-1}
    (\mathbf{x}_{n+1}-\mathbf{x}_n)
\end{equation}

Se considera convergido cuando $d^2 < n \cdot \epsilon_{\text{conv}}$, donde
$\epsilon_{\text{conv}} = 0.1$ en nuestro c\'odigo.  F\'isicamente, esto
significa que el paso es peque\~no comparado con la incertidumbre posterior.
M\'aximo 20 iteraciones (30 para Nivel~2b).


% ======================================================================
\section{Regularizaci\'on de Phillips--Twomey}\label{sec:pt}
% ======================================================================

\subsection{Formulaci\'on}

A diferencia de LM y OE que iteran sobre el modelo no lineal $F(\mathbf{x})$,
el m\'etodo Phillips--Twomey (Phillips, 1962; Twomey, 1963) resuelve una
versi\'on \textbf{linealizada} del problema.  Se linealiza $F$ alrededor de un
punto de referencia (t\'ipicamente la soluci\'on OE):
\begin{equation}
    \mathbf{y} \approx F(\mathbf{x}_0) + \mathbf{K}(\mathbf{x} - \mathbf{x}_0)
    = \mathbf{K}\mathbf{x} + \text{const}
\end{equation}

donde $\mathbf{K}$ es el Jacobiano evaluado en $\mathbf{x}_0$ (en nuestro
c\'odigo se toma de la \'ultima iteraci\'on de OE: \texttt{res\_oe2["K"]}).

El problema linealizado minimiza:
\begin{equation}
    J(\mathbf{x}) = \underbrace{\|\mathbf{y} - \mathbf{K}\mathbf{x}\|^2}_{\text{ajuste}}
    + \underbrace{\gamma\,\|\mathbf{H}\mathbf{x}\|^2}_{\text{penalizaci\'on}}
\end{equation}

donde cada t\'ermino se define como:
\begin{itemize}
    \item $\mathbf{K} \in \mathbb{R}^{m \times n}$: Jacobiano del modelo directo.
    \item $\mathbf{H} \in \mathbb{R}^{p \times n}$: \textbf{matriz de restricci\'on}
          (constraint matrix) que penaliza propiedades no deseadas en la
          soluci\'on (ver abajo).
    \item $\gamma > 0$: par\'ametro de regularizaci\'on (misma funci\'on que en
          Tikhonov).
\end{itemize}

Derivando $J$ respecto a $\mathbf{x}$ e igualando a cero se obtiene la
\textbf{soluci\'on cerrada} (no iterativa):
\begin{equation}\label{eq:pt_solution}
    \boxed{\hat{\mathbf{x}} = (\mathbf{K}^T\mathbf{K}
    + \gamma\,\mathbf{H}^T\mathbf{H})^{-1}\mathbf{K}^T\mathbf{y}}
\end{equation}

\subsection{Matrices de Restricci\'on $\mathbf{H}$}

La elecci\'on de $\mathbf{H}$ determina qu\'e propiedad de la soluci\'on se
penaliza:

\textbf{Orden 0} ($\mathbf{H} = \mathbf{I}_n$, penaliza la \emph{magnitud}):
\begin{equation}
    \|\mathbf{H}\mathbf{x}\|^2 = \|\mathbf{x}\|^2 = \sum_i x_i^2
\end{equation}

Equivalente a Tikhonov cl\'asico con $\mathbf{x}_a = \mathbf{0}$.

\textbf{Orden 1} ($\mathbf{H} = \mathbf{D}_1$, penaliza el \emph{gradiente}):
\begin{equation}
    \mathbf{D}_1 = \begin{bmatrix}
        -1 & 1 & & \\
        & -1 & 1 & \\
        & & \ddots & \ddots
    \end{bmatrix}_{(n-1)\times n}
    \qquad
    \|\mathbf{D}_1\mathbf{x}\|^2 = \sum_i (x_{i+1}-x_i)^2
\end{equation}

Produce soluciones suaves (sin saltos bruscos).

\textbf{Orden 2} ($\mathbf{H} = \mathbf{D}_2$, penaliza la \emph{curvatura}):
\begin{equation}
    \mathbf{D}_2 = \begin{bmatrix}
        1 & -2 & 1 & & \\
        & 1 & -2 & 1 & \\
        & & \ddots & \ddots & \ddots
    \end{bmatrix}_{(n-2)\times n}
    \qquad
    \|\mathbf{D}_2\mathbf{x}\|^2 = \sum_i (x_{i+2}-2x_{i+1}+x_i)^2
\end{equation}

Produce soluciones muy suaves, penalizando cambios en la pendiente.
\textbf{Nuestro c\'odigo usa orden 2} para el an\'alisis L-curve del Nivel~2.

\subsection{Selecci\'on de $\gamma$: M\'etodo de la L-Curve}

El problema cr\'itico en cualquier regularizaci\'on es elegir $\gamma$.
El m\'etodo de la \textbf{L-curve} (Hansen, 1992) es un enfoque gr\'afico:
\begin{enumerate}
    \item Para cada $\gamma$ en un rango (nuestro c\'odigo:
          $\gamma \in [10^{-4}, 10^2]$, 30 puntos logar\'itmicamente
          espaciados), se resuelve~\eqref{eq:pt_solution}.
    \item Se calcula la \textbf{norma del residuo}
          $\rho(\gamma) = \|\mathbf{y}-\mathbf{K}\hat{\mathbf{x}}(\gamma)\|$
          y la \textbf{norma de la soluci\'on}
          $\eta(\gamma) = \|\mathbf{H}\hat{\mathbf{x}}(\gamma)\|$.
    \item Se grafica $(\log\rho, \log\eta)$: la curva resultante tiene forma
          de ``L''.
\end{enumerate}

\begin{equation}
    \text{eje-}x: \log\|\mathbf{y}-\mathbf{K}\hat{\mathbf{x}}\|, \qquad
    \text{eje-}y: \log\|\mathbf{H}\hat{\mathbf{x}}\|
\end{equation}

El $\gamma$ \'optimo est\'a en la \textbf{esquina} de la curva L:
\begin{itemize}
    \item Rama horizontal ($\gamma$ peque\~no): buen ajuste a datos
          ($\rho$ peque\~no) pero soluci\'on irregular ($\eta$ grande).
    \item Rama vertical ($\gamma$ grande): soluci\'on suave ($\eta$ peque\~no)
          pero mal ajuste ($\rho$ grande).
    \item Esquina: balance \'optimo entre ajuste y suavidad.
\end{itemize}

\begin{physics}{Por qu\'e Phillips--Twomey es solo una demostraci\'on}
PT requiere un modelo \emph{lineal} $\mathbf{y} = \mathbf{K}\mathbf{x}$,
pero nuestro modelo DISORT es no lineal.  Por eso usamos el $\mathbf{K}$ de
OE (que ya linealiz\'o alrededor de la soluci\'on).  En la pr\'actica, LM y
OE son superiores para recuperaci\'on porque iteran sobre el modelo no lineal
completo.  PT se usa en nuestro c\'odigo \'unicamente para ilustrar el
concepto de L-curve y la selecci\'on de $\gamma$.
\end{physics}


% ======================================================================
\section{Regularizaci\'on Total Variation (TV)}\label{sec:tv}
% ======================================================================

\subsection{Motivaci\'on}

Los m\'etodos anteriores (Tikhonov, OE, PT) penalizan la norma $L^2$ de las
diferencias, lo que produce soluciones \textbf{suaves} (derivable en todas
partes).  Sin embargo, en perfiles atmosf\'ericos reales pueden existir
\textbf{discontinuidades} (por ejemplo, una capa de aerosol con bordes
abruptos).  La regularizaci\'on $L^2$ suaviza estos bordes, perdiendo
informaci\'on.

La \textbf{Total Variation} (TV) (Rudin et al., 1992) usa la norma $L^1$ en
lugar de $L^2$:
\begin{equation}
    J(\mathbf{x}) = \|\mathbf{y} - \mathbf{K}\mathbf{x}\|^2
    + \gamma\,\underbrace{\text{TV}(\mathbf{x})}_{\text{norma } L^1}
\end{equation}

donde la \textbf{variaci\'on total} discreta es:
\begin{equation}
    \text{TV}(\mathbf{x}) = \sum_{i=1}^{n-1}|x_{i+1} - x_i|
\end{equation}

La diferencia clave con $L^2$ es que $|z|$ penaliza diferencias grandes
\emph{linealmente} (no cuadr\'aticamente), lo que permite que la soluci\'on
tenga saltos bruscos sin ser excesivamente penalizada.

\subsection{M\'etodo IRLS (Iteratively Reweighted Least Squares)}

La norma $L^1$ no es diferenciable en cero, lo que impide el uso directo de
m\'etodos de m\'inimos cuadrados.  El truco es aproximarla iterativamente
como una norma $L^2$ ponderada (\textbf{IRLS}):

En cada iteraci\'on $k$:
\begin{enumerate}
    \item \textbf{Calcular pesos} a partir de la soluci\'on actual:
    \begin{equation}
        w_i^{(k)} = \frac{1}{|x_{i+1}^{(k)} - x_i^{(k)}| + \epsilon}
    \end{equation}
    donde $\epsilon = 10^{-6}$ evita divisi\'on por cero.  Donde las
    diferencias son \emph{grandes} (bordes), $w_i$ es \emph{peque\~no}
    $\to$ menos penalizaci\'on $\to$ el borde se preserva.  Donde las
    diferencias son \emph{peque\~nas} (regiones suaves), $w_i$ es
    \emph{grande} $\to$ m\'as penalizaci\'on $\to$ se suaviza.

    \item \textbf{Ensamblar la penalizaci\'on ponderada}:
    \begin{align}
        \mathbf{W}^{(k)} &= \text{diag}(w_1^{(k)},\ldots,w_{n-1}^{(k)})\\
        \mathbf{H}^{(k)} &= \sqrt{\mathbf{W}^{(k)}}\,\mathbf{D}_1
    \end{align}

    \item \textbf{Resolver el subproblema de Tikhonov ponderado}:
    \begin{equation}
        \mathbf{x}^{(k+1)} = (\mathbf{K}^T\mathbf{K}
        + \gamma\,(\mathbf{H}^{(k)})^T\mathbf{H}^{(k)})^{-1}\mathbf{K}^T\mathbf{y}
    \end{equation}

    \item \textbf{Verificar convergencia}:
    $\|\mathbf{x}^{(k+1)} - \mathbf{x}^{(k)}\| / \|\mathbf{x}^{(k)}\| < 10^{-6}$.
\end{enumerate}

Nuestro c\'odigo usa hasta 20 iteraciones IRLS (t\'ipicamente converge en
3--10).

\begin{concept}{Estado actual: TV no se usa en los niveles 1--3}
TV est\'a implementado en \texttt{TotalVariation} pero no se utiliza en
ninguno de los tres niveles porque:
\begin{itemize}
    \item Los perfiles de aerosol que modelamos son suaves (distribuci\'on
          exponencial), sin discontinuidades.
    \item Con solo $n = 1$ o $n = 3$ par\'ametros, la regularizaci\'on TV no
          aporta ventaja sobre Tikhonov/OE.
    \item TV ser\'ia relevante para recuperar perfiles verticales de AOD con
          bordes abruptos (por ejemplo, capas de polvo sahariano).
\end{itemize}
\end{concept}


% ======================================================================
\part{Integraci\'on: Del Modelo F\'isico a la Recuperaci\'on de Aerosoles}
% ======================================================================

% ======================================================================
\section{Pipeline Completo: Paso a Paso}
% ======================================================================

A continuaci\'on se describe c\'omo todas las piezas desarrolladas en las
partes anteriores se integran en un flujo de trabajo coherente para resolver
el problema directo-inverso completo.

\subsection{Paso 1: Construcci\'on del Perfil Atmosf\'erico}

\begin{enumerate}
    \item Se inicializa la \textbf{atm\'osfera est\'andar} US-1976
          (\texttt{StandardAtmosphere}) con los perfiles de $T(z)$, $P(z)$,
          $n_{\text{air}}(z)$ calculados por la ecuaci\'on hidrost\'atica
          (Secci\'on~\ref{sec:us76}).

    \item Se calcula el \textbf{espesor \'optico de Rayleigh} por capa
          $\tau_{\text{Ray},k}$ a la longitud de onda deseada
          (Secci\'on~\ref{sec:rayleigh}).

    \item Se computan las \textbf{propiedades \'opticas del aerosol}
          ($\omega_{\text{aer}}$, $g_{\text{aer}}$) a partir de la teor\'ia
          de Mie integrada sobre una distribuci\'on de tama\~nos lognormal
          (Secci\'on~\ref{sec:mie}).

    \item Se distribuye el AOD total verticalmente usando un perfil exponencial
          $\tau_{\text{aer}}(z) \propto e^{-z/H}$ con altura de escala $H$,
          y se combinan Rayleigh + aerosol para obtener
          $(\tau_k, \omega_{0,k}, g_k)$ por capa.

    \item El resultado es un \texttt{AtmosphereProfile} con 20 capas listo
          para el solver.
\end{enumerate}

\subsection{Paso 2: Modelo Directo (Forward)}

\begin{enumerate}
    \item \textbf{Delta-M scaling} (Secci\'on~\ref{sec:deltam}):
          Se escalan $(\tau_k, \omega_{0,k}, g_k) \to (\tau'_k, \omega'_{0,k}, g'_k)$
          para truncar el pico de forward scattering.

    \item \textbf{Eigen-decomposici\'on por capa}: Para cada capa $k$:
    \begin{enumerate}
        \item Se construye la matriz de fase $\mathbf{P}_k$ via
              Ec.~\eqref{eq:phase_matrix}.
        \item Se construye la matriz del sistema
              $\mathbf{A}_k = \mathbf{M}^{-1}(\mathbf{I}_N - \frac{\omega'_{0,k}}{2}\mathbf{P}_k\mathbf{W})$.
        \item Se computan autovalores y autovectores:
              $\mathbf{A}_k\mathbf{E}_k = \mathbf{E}_k\bm{\Lambda}_k$.
        \item Se ordenan los eigenvalores: $\text{Re}(\lambda_j) \geq 0$ primero,
              $\text{Re}(\lambda_j) < 0$ despu\'es (estabilizaci\'on,
              Secci\'on~\ref{sec:scaling}).
        \item Se resuelve el sistema particular
              $(\mathbf{A}_k + \mu_0^{-1}\mathbf{I}_N)\mathbf{Z}_{p,k} = -\mathbf{S}_k$.
    \end{enumerate}

    \item \textbf{Ensamblaje del sistema global de frontera}: Se construye la
          matriz $\mathbf{B}$ de tama\~no $NK \times NK$ con:
    \begin{itemize}
        \item $N/2$ ecuaciones de TOA (Ec.~\ref{eq:bc_toa}),
        \item $N(K-1)$ ecuaciones de continuidad (Ec.~\ref{eq:bc_interface_expanded}),
        \item $N/2$ ecuaciones de BOA con superficie Lambertiana (Ec.~\ref{eq:bc_boa}).
    \end{itemize}
    Todos los t\'erminos exponenciales usan la forma estabilizada
    (Ec.~\ref{eq:stable_exp}).

    \item \textbf{Resoluci\'on}: $\mathbf{B}\hat{\mathbf{c}} = \mathbf{b}$
          mediante \texttt{scipy.linalg.solve}.

    \item \textbf{Extracci\'on de radiancias TOA}: Se eval\'ua la
          Ec.~\eqref{eq:I_toa} para obtener el vector de observaci\'on
          $\mathbf{y} = F(\mathbf{x})$.
\end{enumerate}

\subsection{Paso 3: Generaci\'on de Observaciones Sint\'eticas}

Para validar la inversi\'on, se generan observaciones sint\'eticas:
\begin{equation}
    \mathbf{y}_{\text{obs}} = F(\mathbf{x}_{\text{true}})
    + \bm{\epsilon}, \qquad
    \bm{\epsilon} \sim \mathcal{N}(\mathbf{0}, \mathbf{S}_\epsilon)
\end{equation}

con $\mathbf{S}_\epsilon = \text{diag}(\sigma_1^2,\ldots,\sigma_m^2)$ donde
$\sigma_i = \delta \cdot y_i$ ($\delta = 0.02$ para 2\% de ruido relativo).

\subsection{Paso 4: Inversi\'on (Retrieval)}

Se aplican los m\'etodos de inversi\'on:

\begin{enumerate}
    \item \textbf{LM + Tikhonov} (\texttt{InverseOptimizer}):
          Minimiza la Ec.~\eqref{eq:cost_tikhonov} usando
          \texttt{scipy.optimize.least\_squares} con residuo aumentado.
          R\'apido y robusto para problemas peque\~nos.

    \item \textbf{Optimal Estimation} (\texttt{OptimalEstimation}):
          Itera la Ec.~\eqref{eq:oe_update} hasta convergencia.
          Proporciona diagn\'osticos completos: $\hat{\mathbf{S}}$,
          $\mathbf{A}$, DFS, funci\'on de costo.

    \item \textbf{Phillips--Twomey} (\texttt{PhillipsTwomey}):
          Inversi\'on lineal con la Ec.~\eqref{eq:pt_solution} usando
          $\mathbf{K}$ del \'ultimo paso de OE.  Se usa L-curve para
          seleccionar $\gamma$.

    \item \textbf{Total Variation} (\texttt{TotalVariation}):
          IRLS para preservar bordes en el perfil recuperado.
\end{enumerate}

\subsection{Paso 5: Diagn\'osticos y Validaci\'on}

\begin{enumerate}
    \item \textbf{Self-consistency}: $F(\mathbf{x}_{\text{true}}) \to$
          retrieval $\to$ $|\hat{\mathbf{x}} - \mathbf{x}_{\text{true}}|
          / |\mathbf{x}_{\text{true}}| < 0.01\%$?

    \item \textbf{Averaging kernels}: $\mathbf{A} \approx \mathbf{I}$?
          DFS $\approx n$?

    \item \textbf{Incertidumbre posterior}: $\sqrt{\text{diag}(\hat{\mathbf{S}})}$
          es consistente con el error real del retrieval?

    \item \textbf{Sensibilidad al ruido}: Para m\'ultiples realizaciones de
          ruido, $\text{std}(\hat{\mathbf{x}}) \approx
          \sqrt{\text{diag}(\hat{\mathbf{S}})}$?

    \item \textbf{Sensibilidad al a priori}: Con $\mathbf{S}_a$ grande (a priori
          d\'ebil), el resultado debe ser independiente de $\mathbf{x}_a$.
\end{enumerate}


% ======================================================================
\section{Resumen de la Cadena Matem\'atica}
% ======================================================================

La siguiente cadena resume la transformaci\'on matem\'atica completa desde
los par\'ametros f\'isicos hasta la recuperaci\'on:

\begin{equation*}
    \underbrace{(r_g, \sigma_g, m)}_{\text{Microf\'isica}}
    \xrightarrow{\text{Mie}}
    \underbrace{(\omega_0, g)}_{\text{\'Optica}}
    \xrightarrow{\text{US-76 + Rayleigh}}
    \underbrace{(\tau_k, \omega_{0,k}, g_k)}_{\text{Perfil}}
    \xrightarrow{\delta\text{-M}}
    \underbrace{(\tau'_k, \omega'_{0,k}, g'_k)}_{\text{Escalado}}
\end{equation*}

\begin{equation*}
    \xrightarrow{\text{Eigen}}
    \underbrace{(\lambda_j, E_{ij}, Z_{p,i})_k}_{\text{Modos}}
    \xrightarrow{\text{BC}}
    \underbrace{\hat{C}_{k,j}}_{\text{Constantes}}
    \xrightarrow{\text{Eval.\ TOA}}
    \underbrace{I_i^{\text{TOA}}}_{\text{Radiancias}}
    \xrightarrow{\text{Inversi\'on}}
    \underbrace{\hat{\mathbf{x}}}_{\text{Retrieval}}
\end{equation*}

Cada paso es trazable, verificable y tiene significado f\'isico claro.


% ======================================================================
\section{Resultados de Validaci\'on}
% ======================================================================

\subsection{Forward Model: 4/4 Tests}

\begin{enumerate}
    \item \textbf{L\'imite $\tau \to 0$}: Sin atm\'osfera, la radiancia TOA
          es solo la reflexi\'on directa de superficie:
          $I \approx (A_s/\pi)\mu_0 F_0$. \textbf{PASS}.

    \item \textbf{Convergencia con $N$}: Las radiancias convergen para
          $N = 4 \to 8 \to 16 \to 32$. \textbf{PASS}.

    \item \textbf{Consistencia mono/multi-capa}: Una capa con $\tau=1$ produce
          los mismos resultados que 20 capas con $\tau=0.05$ cada una.
          $\Delta_{\text{rel}} < 10^{-6}$. \textbf{PASS}.

    \item \textbf{Absorci\'on pura}: Con $\omega_0 \approx 0$ y $A_s = 0$,
          la radiancia TOA es m\'inima. \textbf{PASS}.
\end{enumerate}

\subsection{Inverse Model: 3/3 Tests}

\begin{enumerate}
    \item \textbf{Self-consistency (sin ruido)}: Error de retrieval
          $< 0.01\%$. \textbf{PASS}.

    \item \textbf{OE self-consistency}: Retrieval converge con DFS$\approx 1$.
          \textbf{PASS}.

    \item \textbf{Sensibilidad al a priori}: Con $\mathbf{S}_a$ grande,
          el resultado es independiente de $\mathbf{x}_a$ (spread $< 10^{-4}$).
          \textbf{PASS}.
\end{enumerate}


% ======================================================================
\section{Par\'ametros del Modelo}
% ======================================================================

\begin{table}[H]
\centering
\small
\begin{tabular}{@{}llp{6.5cm}l@{}}
\toprule
\textbf{Par\'ametro} & \textbf{S\'imbolo} & \textbf{Descripci\'on} & \textbf{Rango} \\ \midrule
Espesor \'Optico & $\tau$ (AOD) & Extinci\'on total integrada verticalmente & $[0.01, 5]$ \\
Albedo Disp.\ Simple & $\omega_0$ & Fracci\'on dispersi\'on/extinci\'on & $(0, 1)$ \\
Par\'am.\ Asimetr\'ia & $g$ & Preferencia angular de dispersi\'on & $(-1, 1)$ \\
Albedo Superficie & $A_s$ & Reflectividad Lambertiana & $[0, 1]$ \\
\'Ang.\ Cenital Solar & $\theta_0$ & Posici\'on del sol & $[0, 85\degree]$ \\
Flujo Solar & $F_0$ & Irradiancia TOA ($\pi$ por convenci\'on) & W\,m$^{-2}$ \\
N\'um.\ Streams & $N$ & Direcciones angulares discretas & $\{4,8,16,32\}$ \\
Longitud de Onda & $\lambda$ & Longitud de onda de c\'alculo & nm \\
Altura Escala Aerosol & $H$ & Distribuci\'on vertical del aerosol & km \\
Regularizaci\'on & $\gamma$ & Balance datos/suavidad & $> 0$ \\
\bottomrule
\end{tabular}
\caption{Par\'ametros del modelo y su significado f\'isico.}
\end{table}

\begin{table}[H]
\centering
\small
\begin{tabular}{@{}llp{6.5cm}l@{}}
\toprule
\textbf{Par\'ametro} & \textbf{S\'imbolo} & \textbf{Descripci\'on} & \textbf{Valor} \\ \midrule
Exp.\ \AA ngstr\"om & $\alpha$ & Dependencia espectral del AOD & $\sim 1.5$ \\
Fracci\'on Difusa & $f_d$ & DHI/GHI observable SURFRAD & $[0, 1]$ \\
Bandas BB & $\lambda_b$ & Longitudes de onda del modelo 3-bandas & 400, 550, 800\,nm \\
\bottomrule
\end{tabular}
\caption{Par\'ametros adicionales del Nivel 3 (SURFRAD).}
\end{table}

\begin{table}[H]
\centering
\small
\begin{tabular}{@{}lcccc@{}}
\toprule
\textbf{Tipo Aerosol} & $r_g$ ($\mu$m) & $\sigma_g$ & $n_r$ & $n_i$ \\ \midrule
Continental     & 0.03 & 2.24 & 1.53 & 0.008 \\
Urbano          & 0.03 & 2.24 & 1.55 & 0.025 \\
Mar\'itimo       & 0.40 & 2.51 & 1.38 & 0.001 \\
Polvo Des\'ertico & 0.50 & 2.20 & 1.53 & 0.005 \\
Quema Biomasa   & 0.07 & 1.80 & 1.52 & 0.020 \\
\bottomrule
\end{tabular}
\caption{Presets de tipos de aerosol (basados en OPAC y Dubovik).}
\end{table}


% ======================================================================
\section{Implementaci\'on del Nivel 2: Atm\'osfera de 20 Capas}\label{sec:level2_impl}
% ======================================================================

El Nivel~2 representa la configuraci\'on realista del modelo: una atm\'osfera
multicapa con dispersi\'on de Rayleigh y aerosoles distribuidos verticalmente,
donde se recuperan simult\'aneamente tres par\'ametros:
$\mathbf{x} = [\text{AOD},\,\omega_0,\,g]^T$.

\subsection{Configuraci\'on del Aerosol}

Se utiliza el preset \textbf{continental} de la Tabla~3, con los par\'ametros
microf\'isicos:
\begin{center}
    $r_g = 0.03\;\mu$m, \quad $\sigma_g = 2.24$, \quad
    $m = 1.53 + 0.008i$
\end{center}

A la longitud de onda de referencia $\lambda = 550$\,nm, la teor\'ia de Mie
(Secci\'on~\ref{sec:mie}) integrada sobre la distribuci\'on lognormal
(200 radios en $[r_g/10,\, r_g \cdot 10]$ con integraci\'on trapezoidal)
produce las propiedades \'opticas \emph{bulk}:
\begin{center}
    $\omega_{\text{aer}} = \text{SSA}_{\text{Mie}}$, \qquad
    $g_{\text{aer}} = g_{\text{Mie}}$
\end{center}

Estos valores dependen \'unicamente de la microf\'isica y la longitud de onda;
no del AOD, que es un par\'ametro libre de la inversi\'on.

\subsection{Construcci\'on del Perfil Atmosf\'erico}

La cadena de construcci\'on es:
\begin{enumerate}
    \item \texttt{StandardAtmosphere($\lambda$)}: genera los 20 niveles de la
          US-1976 (Secci\'on~\ref{sec:us76}) y calcula $\tau_{\text{Ray},k}$
          por capa.

    \item \texttt{to\_profile(aod\_total, ssa\_aer, g\_aer, H\_aer\_km)}:
    \begin{enumerate}
        \item Distribuye el AOD total con el perfil exponencial
              (Secci\'on~\ref{sec:aer_vertical}):
              \[
                  \tau_{\text{aer},k} = \text{AOD}_{\text{total}}
                  \cdot \frac{e^{-z_{\text{mid},k}/H}}
                  {\sum_{j=1}^{20} e^{-z_{\text{mid},j}/H} \cdot \Delta z_j}
                  \cdot \Delta z_k
              \]
              con $H = 2.0$\,km.

        \item Combina Rayleigh y aerosol capa por capa:
              \begin{align*}
                  \tau_k &= \tau_{\text{Ray},k} + \tau_{\text{aer},k} \\
                  \omega_{0,k} &= \frac{\tau_{\text{Ray},k} + \omega_{\text{aer}}\,\tau_{\text{aer},k}}{\tau_k} \\
                  g_k &= \frac{\omega_{\text{aer}}\,g_{\text{aer}}\,\tau_{\text{aer},k}}{\tau_{\text{Ray},k} + \omega_{\text{aer}}\,\tau_{\text{aer},k}}
              \end{align*}
              (Rayleigh puro: $\omega_{\text{Ray}} = 1$, $g_{\text{Ray}} = 0$.)
    \end{enumerate}

    \item El resultado es un \texttt{AtmosphereProfile} con 20 capas, listo
          para el solver DISORT con $N = 36$ streams.
\end{enumerate}

\subsection{Modelo Directo y Observaciones}

El solver DISORT resuelve el problema directo con delta-M scaling activado,
produciendo $m = N/2 = 18$ radiancias TOA upwelling (en los nodos de
cuadratura ascendentes). Adicionalmente, se calculan radiancias en \'angulos
arbitrarios mediante la integraci\'on de funci\'on fuente
(Ecs.~35a,b del reporte).

Las observaciones sint\'eticas se generan como:
\begin{equation}
    \mathbf{y}_{\text{obs}} = F(\mathbf{x}_{\text{true}}) + \bm{\epsilon},
    \qquad \bm{\epsilon} \sim \mathcal{N}(\mathbf{0},\,\mathbf{S}_\epsilon)
\end{equation}
con ruido relativo configurable (t\'ipicamente $\delta = 0\%$ para
validaci\'on, $2\%$ para pruebas de robustez).

\subsection{Configuraci\'on de la Inversi\'on}

Se aplican dos m\'etodos principales:

\paragraph{LM + Tikhonov}
\begin{itemize}
    \item A priori: $\mathbf{x}_a = [0.5,\; 0.85,\; 0.5]^T$
    \item Regularizaci\'on: $\gamma = 0.05$
    \item L\'imites f\'isicos: $\text{AOD}\in[0.01, 3.0]$,
          $\omega_0\in[0.5, 0.9999]$, $g\in[0.0, 0.999]$
    \item Guess inicial: $\mathbf{x}_0 = [0.2,\; 0.90,\; 0.60]^T$
\end{itemize}

\paragraph{Optimal Estimation}
\begin{itemize}
    \item A priori: $\mathbf{x}_a = [0.5,\; 0.85,\; 0.5]^T$
    \item Covarianza a priori: $\mathbf{S}_a = \text{diag}(0.3^2,\; 0.1^2,\; 0.2^2)
          = \text{diag}(0.09,\; 0.01,\; 0.04)$
    \item Covarianza del error: $\mathbf{S}_\epsilon = (\sigma)^2\,\mathbf{I}_m$
          con $\sigma = 0.02\,\bar{y}$
    \item M\'aximo 20 iteraciones, convergencia Mahalanobis $< 0.1$
    \item Amortiguamiento LM: $\gamma_{\text{LM}} = 10$
\end{itemize}

\paragraph{Phillips--Twomey (demostraci\'on)}
Se ejecuta la L-curve con $\gamma \in [10^{-4}, 10^2]$ (30 puntos
logar\'itmicamente espaciados) usando el Jacobiano $\mathbf{K}$ de la
\'ultima iteraci\'on de OE, con restricci\'on de orden 2.

\subsection{Extensi\'on Multi-longitud de Onda}

El modelo tambi\'en genera radiancias TOA a m\'ultiples longitudes de onda
$\lambda \in \{440, 550, 670, 870\}$\,nm para caracterizar la dependencia
espectral. Para cada $\lambda$:
\begin{enumerate}
    \item Se recalculan las propiedades de Mie: $\omega_{\text{aer}}(\lambda)$,
          $g_{\text{aer}}(\lambda)$.
    \item Se recalcula $\tau_{\text{Ray},k}(\lambda) \propto \lambda^{-4.04}$.
    \item Se mantiene el mismo AOD total (referencia a 550\,nm).
    \item Se resuelve el problema directo independientemente.
\end{enumerate}

Esto permite visualizar c\'omo la dispersi\'on de Rayleigh domina en longitudes
de onda cortas ($\tau_{\text{Ray}} \gg \tau_{\text{aer}}$ a 440\,nm) mientras
que el aerosol domina a longitudes de onda largas (870\,nm).

\begin{table}[H]
\centering
\small
\caption{Par\'ametros espec\'ificos del Nivel 2.}\label{tab:level2_params}
\begin{tabular}{@{}lll@{}}
\toprule
\textbf{Par\'ametro} & \textbf{S\'imbolo/Variable} & \textbf{Valor} \\ \midrule
Tipo de aerosol & \texttt{L2\_AEROSOL\_TYPE} & continental \\
Longitud de onda & $\lambda$ & 550\,nm \\
AOD total & $\tau_{\text{total}}$ & 0.3 \\
Altura de escala & $H_{\text{aer}}$ & 2.0\,km \\
N\'um.\ streams & $N$ & 36 \\
N\'um.\ capas & $K$ & 20 \\
\'Angulo cenital solar & $\theta_0$ & 30° \\
Albedo de superficie & $A_s$ & 0.1 \\
Flujo solar & $F_0$ & $\pi$ \\
\midrule
\multicolumn{3}{@{}l}{\textbf{Inversi\'on LM + Tikhonov}} \\
A priori & $\mathbf{x}_a$ & $[0.5,\; 0.85,\; 0.5]^T$ \\
Regularizaci\'on & $\gamma$ & 0.05 \\
L\'imites inf. & --- & $[0.01,\; 0.5,\; 0.0]$ \\
L\'imites sup. & --- & $[3.0,\; 0.9999,\; 0.999]$ \\
\midrule
\multicolumn{3}{@{}l}{\textbf{Inversi\'on OE}} \\
A priori & $\mathbf{x}_a$ & $[0.5,\; 0.85,\; 0.5]^T$ \\
Covarianza a priori & $\mathbf{S}_a$ & diag$(0.09,\; 0.01,\; 0.04)$ \\
Ruido relativo & $\delta$ & 0.02 (2\%) \\
Amortiguamiento LM & $\gamma_{\text{LM}}$ & 10 \\
\midrule
\multicolumn{3}{@{}l}{\textbf{Phillips--Twomey}} \\
Orden de restricci\'on & --- & 2 (segunda derivada) \\
Rango L-curve & $\gamma$ & $[10^{-4},\; 10^{2}]$ (30 pts) \\
\bottomrule
\end{tabular}
\end{table}


% ======================================================================
\section{Estudio de Recuperaci\'on de 3 Par\'ametros}\label{sec:retrieval_study}
% ======================================================================

Para evaluar la capacidad del sistema de inversi\'on, se realiz\'o un estudio
sistem\'atico de recuperaci\'on simult\'anea de los tres par\'ametros
$\mathbf{x} = [\text{AOD},\,\omega_0,\,g]^T$ usando Optimal Estimation.

\subsection{Dise\~no Experimental}

Se definieron 20 combinaciones de par\'ametros verdaderos
$(\tau_{\text{true}}, \omega_{0,\text{true}}, g_{\text{true}})$ que cubren
el rango realista de aerosoles atmosf\'ericos:
\begin{itemize}
    \item AOD: $[0.05,\, 1.50]$ --- desde atm\'osfera limpia hasta tormenta de polvo.
    \item SSA: $[0.78,\, 0.99]$ --- desde aerosol absorbente hasta casi conservativo.
    \item $g$: $[0.55,\, 0.76]$ --- desde part\'iculas peque\~nas hasta polvo grueso.
\end{itemize}

Para cada caso $i$:
\begin{enumerate}
    \item Se genera la observaci\'on sint\'etica: $\mathbf{y}_i = F(\mathbf{x}_{i,\text{true}})$.
    \item Se configura OE con a priori $\mathbf{x}_a = [0.5,\, 0.88,\, 0.65]^T$
          y covarianzas $\mathbf{S}_a = \text{diag}(0.25, 0.0225, 0.0625)$,
          $\mathbf{S}_\epsilon = \sigma^2 \mathbf{I}$ con $\sigma = 0.02\,\bar{y}$.
    \item Se ejecuta la recuperaci\'on con amortiguamiento LM $\gamma = 0.05$.
\end{enumerate}

\subsection{Resultados}

Los 20 casos convergieron exitosamente. Las m\'etricas de calidad son:
\begin{center}
\begin{tabular}{@{}lccc@{}}
\toprule
\textbf{Par\'ametro} & \textbf{RMSE} & \textbf{Sesgo} & \textbf{$R^2$ impl\'icito}\\ \midrule
AOD ($\tau$) & 0.089 & $-0.020$ & Alto \\
SSA ($\omega_0$) & 0.064 & $-0.023$ & Moderado \\
$g$ (asimetr\'ia) & 0.043 & $-0.012$ & Moderado \\
\bottomrule
\end{tabular}
\end{center}

Los gr\'aficos de dispersi\'on (verdadero vs.\ recuperado) muestran que los puntos
se agrupan cerca de la l\'inea 1:1, con mayor dispersi\'on en SSA y $g$ que en AOD.
Esto es consistente con la teor\'ia: la radiancia TOA tiene mayor sensibilidad
al espesor \'optico total que a los par\'ametros de dispersi\'on simple.


% ======================================================================
\section{Casos F\'isicos Extremos}\label{sec:extreme}
% ======================================================================

Para verificar el comportamiento del solver en condiciones l\'imite, se
analizaron cuatro escenarios con $\text{AOD} = 0.5$ y $\theta_0 = 30°$:

\begin{center}
\small
\begin{tabular}{@{}clccll@{}}
\toprule
\textbf{Panel} & \textbf{Caso} & $\omega_0$ & $A_s$ &
  \textbf{$R$ (TOA)} & \textbf{F\'isica esperada} \\ \midrule
(a) & Sin absorci\'on  & 1.0 & 1.0 & 1.0000 & Toda la energ\'ia escapa \\
(b) & Absorci\'on total & 0.0 & 0.0 & 0.0000 & Nada escapa \\
(c) & Conservativo + negro & 1.0 & 0.0 & 0.0616 & Solo difusa escapa \\
(d) & Absorbente + espejo & 0.0 & 1.0 & 0.2488 & Solo haz directo refleja \\
\bottomrule
\end{tabular}
\end{center}

\begin{keyresult}{Verificaci\'on de L\'imites F\'isicos}
\begin{itemize}
    \item Caso~(a): $R = 1$ confirma conservaci\'on perfecta de energ\'ia ---
          sin absorci\'on en ning\'un lugar, todo el flujo incidente retorna al espacio.
    \item Caso~(b): $R = 0$ confirma que la combinaci\'on $\omega_0 = 0$ (atm\'osfera
          puramente absorbente) y $A_s = 0$ (superficie negra) elimina toda radiaci\'on.
    \item Caso~(c): Con $\omega_0 = 1$, $F_{\text{abs,atm}} = 0$ exactamente,
          verificando que la atm\'osfera no absorbe energ\'ia bajo dispersi\'on conservativa.
    \item Caso~(d): $R = 0.249$ porque el haz directo se aten\'ua a
          $e^{-\tau/\mu_0} = e^{-0.5/0.866} = 0.561$ al descender, y luego se aten\'ua
          nuevamente al ascender despu\'es de reflejarse en la superficie.
\end{itemize}
\end{keyresult}


% ======================================================================
\section{Balance de Energ\'ia}\label{sec:energy_balance}
% ======================================================================

El balance de energ\'ia para el sistema atm\'osfera--superficie se descompone como:
\begin{equation}\label{eq:energy_balance}
    \boxed{F_{\text{in}} = F_\uparrow^{\text{TOA}} + F_{\text{abs,sfc}} + F_{\text{abs,atm}}}
\end{equation}

donde:
\begin{align}
    F_{\text{in}} &= F_0\,\mu_0 & &\text{(flujo solar incidente)} \\
    F_\uparrow^{\text{TOA}} &= 2\pi\sum_i w_i\,\mu_i\,|I_i^{\text{TOA}}|
        & &\text{(flujo que escapa al espacio)} \\
    F_{\text{abs,sfc}} &= (1-A_s)\left(F_{\downarrow}^{\text{diff}} + F_{\text{dir}}\right)
        & &\text{(absorbido por la superficie)} \\
    F_{\text{abs,atm}} &= F_{\text{in}} - F_\uparrow^{\text{TOA}} - F_{\text{abs,sfc}}
        & &\text{(absorbido por la atm\'osfera)}
\end{align}

Se verific\'o esta descomposici\'on en 13 escenarios (barrido de SSA a albedo fijo,
barrido de albedo a SSA fijo, m\'as los 4 casos extremos). El gr\'afico de barras
apiladas muestra que la suma de los tres componentes iguala $F_{\text{in}}$ en todos
los casos, con error m\'aximo $|F_{\text{total}} - F_{\text{in}}| = 0$.


% ======================================================================
\section{Comparaci\'on con Datos AERONET}\label{sec:aeronet}
% ======================================================================

Para validar el modelo contra datos reales, se descargaron observaciones de la
red AERONET (AErosol RObotic NETwork, Holben et al., 1998) del sitio GSFC
(Goddard Space Flight Center, Maryland, EE.UU.) para el per\'iodo
junio--agosto 2023.

\subsection{Datos Utilizados}

\begin{itemize}
    \item \textbf{AOD Level 2.0} (promedios diarios): 87 registros con AOD
          a 440\,nm en el rango $[0.057,\, 2.193]$, media $\bar{\tau} = 0.446$.
    \item \textbf{Inversi\'on SSA} (Almucantar Level 2.0): 25 d\'ias con SSA v\'alido
          a 440\,nm en $[0.920,\, 0.966]$.
    \item \textbf{Inversi\'on Asymmetry Factor-Total}: par\'ametro de asimetr\'ia $g$
          a 440\,nm en $[0.661,\, 0.743]$.
\end{itemize}

\subsection{Metodolog\'ia de Simulaci\'on}

Para cada d\'ia con datos de AOD v\'alidos:
\begin{enumerate}
    \item Se construye un perfil de 1 capa: \texttt{AtmosphereProfile(aod=$\tau_{440}$,
          ssa=$\omega_{0,440}$, g=$g_{440}$, albedo=$A_s$, theta0=$\theta_0$)}.
    \item Si el d\'ia tiene datos de inversi\'on AERONET (SSA y $g$), se usan esos valores;
          si no, se usan valores por defecto ($\omega_0 = 0.92$, $g = 0.70$).
    \item Se ejecuta el solver DISORT y se computa la reflectancia TOA:
          $R = F_\uparrow^{\text{TOA}} / (F_0\,\mu_0)$.
\end{enumerate}

\subsection{Resultados}

Se simularon 87 d\'ias con \'exito. Los gr\'aficos de comparaci\'on muestran:
\begin{itemize}
    \item \textbf{AOD vs.\ $R$ simulado}: Relaci\'on positiva creciente ---
          a mayor carga de aerosol, mayor reflectancia TOA. Los puntos coloreados
          por SSA muestran que aerosoles m\'as absorbentes (SSA bajo) producen
          menor $R$ a igual AOD.
    \item \textbf{Serie temporal}: Las variaciones de AOD y $R$ simulado
          est\'an fuertemente correlacionadas, con eventos de alta carga
          (posiblemente transporte de humo de incendios canadienses en junio 2023)
          claramente visibles.
\end{itemize}


% ======================================================================
\section{Nivel 3: Recuperaci\'on de AOD con Datos Reales SURFRAD}
\label{sec:surfrad}
% ======================================================================

El Nivel~3 cierra el ciclo problema directo $\to$ inverso utilizando
\textbf{datos observacionales reales} para la recuperaci\'on del AOD.  Se
combinan dos redes de medici\'on:
\begin{itemize}
    \item \textbf{AERONET} (Holben et al., 1998): proporciona los par\'ametros
          \'opticos del aerosol (AOD, SSA, $g$) como verdad de referencia.
    \item \textbf{SURFRAD} (Augustine et al., 2000): proporciona flujos de
          radiaci\'on de banda ancha medidos en superficie.
\end{itemize}

\subsection{Datos SURFRAD}

La red SURFRAD (Surface Radiation Budget Network) opera estaciones en EE.UU.\
que miden flujos radiativos de banda ancha con piranometros y pirheliometros:
\begin{itemize}
    \item $\text{GHI}$ --- Global Horizontal Irradiance [W\,m$^{-2}$]:
          flujo total hemisf\'erico descendente.
    \item $\text{DNI}$ --- Direct Normal Irradiance [W\,m$^{-2}$]:
          flujo del haz directo normal al sol.
    \item $\text{DHI}$ --- Diffuse Horizontal Irradiance [W\,m$^{-2}$]:
          flujo difuso hemisf\'erico ($\text{GHI} - \text{DNI}\cos\theta_0$).
\end{itemize}

\begin{physics}{SURFRAD mide flujos, no radiancias angulares}
Es importante notar que los piranometros integran la radiaci\'on sobre todo el
hemisferio (con ponderaci\'on coseno), produciendo un \'unico valor escalar de
flujo.  No proporcionan informaci\'on sobre la distribuci\'on angular de la
radiancia.  Para visualizar campos de radiancia angulares, se simulan con DISORT
usando par\'ametros AERONET reales.
\end{physics}

\subsection{Observable: Fracci\'on Difusa}

El observable para la recuperaci\'on es la \textbf{fracci\'on difusa}:
\begin{equation}\label{eq:diff_frac}
    \boxed{f_d = \frac{\text{DHI}}{\text{GHI}}}
\end{equation}

$f_d$ tiene la ventaja de ser un cociente adimensional que es robusto frente a
errores de calibraci\'on absoluta y depende fuertemente del AOD:
\begin{itemize}
    \item AOD bajo $\to$ $f_d$ peque\~na (la mayor\'ia de la radiaci\'on es directa).
    \item AOD alto $\to$ $f_d$ grande (la radiaci\'on directa se aten\'ua y la
          difusa aumenta por dispersi\'on m\'ultiple).
\end{itemize}

La relaci\'on $f_d(\tau)$ es \textbf{mon\'otona creciente} para par\'ametros
de dispersi\'on fijos ($\omega_0$, $g$), lo que garantiza unicidad de la
soluci\'on y permite el uso de bisecci\'on.

\subsection{Modelo de Banda Ancha: 3 Longitudes de Onda}

Los instrumentos SURFRAD miden la radiaci\'on integrada sobre un rango espectral
amplio (aproximadamente 300--3000\,nm).  Para modelar este flujo de banda ancha
con nuestro solver monocrom\'atico, se utiliza un \textbf{modelo de 3 bandas}:

\begin{center}
\begin{tabular}{@{}cccc@{}}
\toprule
\textbf{Banda} & $\lambda_c$ (nm) & $F_{\text{solar}}$ (W\,m$^{-2}$\,nm$^{-1}$)
& $\Delta\lambda$ (nm) \\ \midrule
VIS  & 400 & 1.60 & 150 \\
VIS  & 550 & 1.88 & 200 \\
NIR  & 800 & 1.14 & 350 \\
\bottomrule
\end{tabular}
\end{center}

El AOD se escala espectralmente usando el exponente de \AA ngstr\"om
$\alpha$ (derivado de AERONET):
\begin{equation}\label{eq:angstrom}
    \tau(\lambda) = \tau_{550}\left(\frac{\lambda}{550}\right)^{-\alpha}
\end{equation}

Para cada banda, se ejecuta DISORT por separado y se combinan los flujos
ponderados por la irradiancia solar y el ancho de banda:
\begin{align}
    F_{\text{dir,BB}} &= \sum_{b=1}^{3} F_{\text{solar},b}\,\Delta\lambda_b\,
    \cos\theta_0\,e^{-\tau_b/\cos\theta_0}\\
    F_{\text{diff,BB}} &= \sum_{b=1}^{3} F_{\text{solar},b}\,\Delta\lambda_b\,
    F_{\downarrow,\text{diff}}^{(b)}
\end{align}

donde $F_{\downarrow,\text{diff}}^{(b)}$ es el flujo difuso descendente
calculado por DISORT para la banda $b$.

\subsection{Algoritmo de Recuperaci\'on: Bisecci\'on}

Dado que $f_d(\tau)$ es mon\'otona, se utiliza el m\'etodo de bisecci\'on para
encontrar el AOD que reproduce la fracci\'on difusa observada:

\begin{enumerate}
    \item Se calcula $f_d^{\text{obs}} = \text{DHI}^{\text{obs}} / \text{GHI}^{\text{obs}}$.
    \item Se define el intervalo de b\'usqueda $\tau \in [0.001, 5.0]$.
    \item En cada iteraci\'on, se eval\'ua $f_d^{\text{model}}(\tau_{\text{mid}})$
          usando el modelo de 3 bandas.
    \item Se actualiza el intervalo seg\'un el signo de
          $f_d^{\text{model}} - f_d^{\text{obs}}$.
    \item Se itera hasta $|\tau_{\text{hi}} - \tau_{\text{lo}}| < 10^{-4}$.
\end{enumerate}

\begin{concept}{Ventajas de la bisecci\'on sobre LM/OE}
\begin{itemize}
    \item \textbf{Sin regularizaci\'on}: no hay sesgo hacia un a priori
          $\mathbf{x}_a$, eliminando el bias sistem\'atico observado en el
          Nivel~1 con LM+Tikhonov.
    \item \textbf{Convergencia garantizada}: la monoton\'ia de $f_d(\tau)$
          asegura convergencia a la soluci\'on \'unica.
    \item \textbf{1D}: con un solo par\'ametro a recuperar (AOD), la
          bisecci\'on es m\'as eficiente que m\'etodos iterativos generales.
\end{itemize}
\end{concept}

\subsection{Colocaci\'on AERONET--SURFRAD}

Para la validaci\'on, se utiliza la estaci\'on SURFRAD de \textbf{Bondville}
(Illinois, 40.05°N, 88.37°W), colocada con el sitio AERONET
\textbf{BONDVILLE}.  Los datos se filtran para condiciones de cielo despejado:
\begin{itemize}
    \item QC flag = 0 (calidad \'optima)
    \item DNI $> 100$ W\,m$^{-2}$ (cielo despejado, sin nubes gruesas)
    \item Horario de mediod\'ia ($\pm 3$ horas del mediod\'ia solar)
    \item Colocalidad temporal: datos de AERONET del mismo d\'ia proporcionan
          los valores de SSA y $g$ necesarios para el modelo
\end{itemize}

\subsection{Resultados del Nivel 3}

La recuperaci\'on de AOD por fracci\'on difusa produce:
\begin{itemize}
    \item $R^2 \approx 0.99$ contra el AOD de AERONET (verdad de referencia)
    \item Sesgo bajo, sin la saturaci\'on observada en inversiones de radiancia
    \item El alto $R^2$ se explica porque:
    \begin{enumerate}
        \item La bisecci\'on no tiene sesgo de regularizaci\'on.
        \item La fracci\'on difusa es mon\'otona en AOD (mapeo 1:1).
        \item Los par\'ametros de dispersi\'on (SSA, $g$) se toman de AERONET,
              eliminando la degeneraci\'on AOD--SSA que limita la precisi\'on
              en inversiones de radiancia angular (Secci\'on~\ref{sec:interp_coupling}).
    \end{enumerate}
\end{itemize}

\subsection{Simulaci\'on del Campo de Radiancia con Datos Reales}

Para visualizar c\'omo se ver\'ia el campo de radiancia angular bajo condiciones
atmosf\'ericas reales (algo que SURFRAD no mide), se seleccionan 5 d\'ias
representativos del dataset AERONET que cubren el rango de AOD observado
(m\'inimo, percentil 25, mediana, percentil 75, m\'aximo) y se simula con
DISORT usando la interpolaci\'on continua (Ecs.~35a,b) en una malla fina de
$1°$ a $89°$:
\begin{equation}
    I^{\text{sim}}(\theta) = I^{\text{DISORT}}(\cos\theta;\,
    \tau^{\text{AERONET}}, \omega_0^{\text{AERONET}}, g^{\text{AERONET}})
    \qquad \theta \in [1°, 89°]
\end{equation}

Esto muestra la transici\'on gradual de campos de radiancia suaves (AOD bajo)
a campos m\'as isotr\'opicos (AOD alto) debido al aumento de la dispersi\'on
m\'ultiple.


% ======================================================================
\section{Comparaci\'on de M\'etodos: LM+Tikhonov vs.\ Optimal Estimation}
\label{sec:lm_vs_oe}
% ======================================================================

Se implement\'o una comparaci\'on sistem\'atica entre los dos m\'etodos
iterativos principales:

\subsection{Diferencias Fundamentales}

\begin{center}
\small
\begin{tabular}{@{}lp{5.5cm}p{5.5cm}@{}}
\toprule
\textbf{Aspecto} & \textbf{LM + Tikhonov} & \textbf{Optimal Estimation} \\ \midrule
Regularizaci\'on & $\gamma\|\mathbf{x}-\mathbf{x}_a\|^2$ (escalar) &
    $(\mathbf{x}-\mathbf{x}_a)^T\mathbf{S}_a^{-1}(\mathbf{x}-\mathbf{x}_a)$
    (matricial) \\
Ponderaci\'on datos & Uniforme ($\|\mathbf{r}\|^2$) &
    $\mathbf{S}_\epsilon^{-1}$ (covarianza de error) \\
Diagn\'osticos & Solo $\mathbf{x}$ y costo & $\hat{\mathbf{S}}$,
    $\mathbf{A}$, DFS, $\mathbf{G}$, $\mathbf{K}$ \\
Implementaci\'on & \texttt{scipy.optimize.least\_squares} &
    Gauss--Newton propio \\
\bottomrule
\end{tabular}
\end{center}

\subsection{Resultados Comparativos}

En los niveles 1 y 2, OE generalmente produce recuperaciones m\'as precisas
que LM+Tikhonov, especialmente para AOD extremos, porque:
\begin{enumerate}
    \item La ponderaci\'on por $\mathbf{S}_\epsilon^{-1}$ da m\'as peso a
          mediciones con menor error.
    \item La covarianza a priori $\mathbf{S}_a$ permite especificar
          incertidumbres \emph{por par\'ametro}, en lugar de un \'unico
          $\gamma$ global.
    \item El DFS indica cu\'anta informaci\'on proviene de los datos vs.\ del
          a priori, guiando la interpretaci\'on del resultado.
\end{enumerate}


% ======================================================================
\part{Interpretaci\'on de Resultados}
% ======================================================================

% ======================================================================
\section{M\'aximos de Radiancia TOA cerca de VZA $\approx 90°$}
\label{sec:interp_vza}
% ======================================================================

\subsection{Origen Matem\'atico}

La radiancia TOA ascendente viene dada por la Ec.~\eqref{eq:I_toa}:
\[
I_i^{\text{TOA}} = \sum_{j=1}^{N}\hat{C}_{1,j}\,E^{(1)}_{ij}\,f_j(0)
+ Z_{p,i}^{(1)}
\]

El t\'ermino dominante para streams con $\mu_i \approx 0$ (VZA $\to 90°$) proviene
de la \textbf{soluci\'on particular} y la estructura de la matriz $\mathbf{A}$:

\begin{equation}\label{eq:A_mu_dependence}
    A_{ii} = \frac{1}{\mu_i}\left(1 - \frac{\omega_0}{2}\,w_i\,P_{ii}\right)
\end{equation}

Cuando $|\mu_i| \to 0$, $A_{ii} \to \pm\infty$. Esto tiene una interpretaci\'on
f\'isica directa.

\subsection{Interpretaci\'on F\'isica: El Efecto del Camino \'Optico}

Un fot\'on que se propaga con \'angulo cenital $\theta$ recorre un camino \'optico
\emph{oblicuo} por cada incremento vertical $d\tau$:
\begin{equation}
    ds = \frac{d\tau}{\mu} = \frac{d\tau}{\cos\theta}
\end{equation}

Para VZA $\to 90°$, $\mu \to 0$ y $ds \to \infty$: el fot\'on atraviesa un
camino \'optico \textbf{infinitamente largo} por unidad de $\Delta\tau$ vertical.
Esto tiene dos consecuencias:

\begin{enumerate}
    \item \textbf{Mayor probabilidad de interacci\'on}: Cada fot\'on dispersado
          hacia \'angulos rasantes ha interactuado muchas veces con el medio,
          acumulando contribuciones de dispersi\'on m\'ultiple.

    \item \textbf{Convergencia de contribuciones}: La radiancia integra sobre
          todas las fuentes de scattering a lo largo del camino oblicuo. Para
          \'angulos rasantes, este camino cruza pr\'acticamente toda la atm\'osfera
          horizontalmente, recogiendo contribuciones de un volumen grande.
\end{enumerate}

Matem\'aticamente, la soluci\'on particular satisface:
\begin{equation}
    \left(\mathbf{A} + \frac{1}{\mu_0}\,\mathbf{I}\right)\mathbf{Z}_p
    = -\mathbf{M}^{-1}\mathbf{Q}
\end{equation}

Para la fila $i$ con $\mu_i$ peque\~no:
\begin{equation}
    Z_{p,i} \propto \frac{Q_i / \mu_i}{1/\mu_i + 1/\mu_0}
    = \frac{Q_i}{\mu_i} \cdot \frac{\mu_0}{\mu_0 + \mu_i}
    \xrightarrow{\mu_i \to 0} \frac{Q_i}{\mu_i}
\end{equation}

El t\'ermino $Q_i/\mu_i$ \textbf{diverge} para $\mu_i \to 0$, lo que explica el
crecimiento de la radiancia hacia \'angulos rasantes. F\'isicamente, esto se
interpreta como la \emph{masa de aire} (air mass factor) $1/\mu$ que amplifica
la fuente de scattering.

\begin{physics}{An\'alogo con la Ley de Beer--Lambert}
La ley de Beer--Lambert para un camino oblicuo es
$I = I_0\,e^{-\tau/\mu}$.  Cuando $\mu \to 0$, la \emph{atenuaci\'on}
tambi\'en se hace infinita.  La competencia entre la mayor fuente de scattering
($\propto 1/\mu$) y la mayor atenuaci\'on ($\propto e^{-\tau/\mu}$) produce un
m\'aximo en la radiancia TOA a VZA alto pero finito (t\'ipicamente $70°$--$85°$),
seguido de un decaimiento cuando la atenuaci\'on domina.
\end{physics}


% ======================================================================
\section{No Linealidad en la Recuperaci\'on del AOD (Nivel 1)}
\label{sec:interp_nonlinear}
% ======================================================================

\subsection{El Modelo Directo es Intr\'insecamente No Lineal}

La relaci\'on entre AOD y radiancia TOA no es lineal. Esto se demuestra
expandiendo la dependencia funcional:

\begin{equation}
    I^{\text{TOA}}(\tau) = \sum_j C_j(\tau)\,E_{ij}\,e^{-\lambda_j(\tau)\cdot 0}
    + Z_{p,i}(\tau)
\end{equation}

donde $C_j$, $\lambda_j$ y $Z_{p,i}$ \textbf{todos} dependen de $\tau$ de
manera no lineal:

\begin{enumerate}
    \item \textbf{Eigenvalores $\lambda_j(\tau)$}: Provienen de la
          eigendescomposici\'on de $\mathbf{A} = \mathbf{M}^{-1}(\mathbf{I} -
          \frac{\omega_0}{2}\mathbf{P}\mathbf{W})$.  Aunque $\mathbf{A}$ no
          depende expl\'icitamente de $\tau$ dentro de una capa, los eigenvalores
          var\'ian con el delta-M scaling: $\tau' = \tau(1 - \omega_0 f)$, y los
          escalados $\omega_0'$ y $g'$ cambian si se modifica la composici\'on.

    \item \textbf{Constantes de integraci\'on $C_j(\tau)$}: Las condiciones de
          frontera involucran exponenciales $e^{-\lambda_j\Delta\tau}$ y el
          t\'ermino de superficie $e^{-\tau_{\text{total}}/\mu_0}$, ambos
          funciones no lineales de $\tau$.

    \item \textbf{Soluci\'on particular $Z_{p,i}(\tau)$}: Depende del vector
          fuente $\mathbf{Q}$ que escala con $\omega_0$ y $F_0$, y es
          modulado por $e^{-\tau/\mu_0}$ en las condiciones de frontera.
\end{enumerate}

\subsection{Consecuencia para el Problema Inverso}

El Jacobiano $K_{ij} = \partial y_i / \partial x_j$ var\'ia con el punto
de evaluaci\'on $\mathbf{x}$, requiriendo una actualizaci\'on iterativa:

\begin{equation}
    \mathbf{x}_{n+1} = \mathbf{x}_a + \mathbf{G}_n
    \left[\mathbf{y}_{\text{obs}} - F(\mathbf{x}_n) +
    \mathbf{K}_n(\mathbf{x}_n - \mathbf{x}_a)\right]
\end{equation}

Si la relaci\'on fuera lineal ($F(\mathbf{x}) = \mathbf{K}\mathbf{x} + \mathbf{b}$),
una sola iteraci\'on bastar\'ia.  La necesidad de m\'ultiples iteraciones
(t\'ipicamente 3--10) evidencia la no linealidad.

\subsection{Manifestaci\'on en los Gr\'aficos de AOD vs.\ $R$}

En el gr\'afico de dispersi\'on AOD vs.\ reflectancia $R$, se observa:
\begin{itemize}
    \item Para AOD bajo ($\tau < 0.2$): crecimiento aproximadamente lineal.
          La atm\'osfera es \'opticamente delgada y la radiancia escala como
          $I \propto \omega_0\,\tau\,P(\mu,\mu_0)/(4\pi\mu)$.

    \item Para AOD moderado ($0.2 < \tau < 1$): la relaci\'on se curva hacia
          arriba por dispersi\'on m\'ultiple acumulada.

    \item Para AOD alto ($\tau > 1$): \textbf{saturaci\'on} --- la radiancia
          crece cada vez m\'as lentamente porque la atenuaci\'on del haz directo
          ($e^{-\tau/\mu_0}$) domina sobre la ganancia por scattering.
\end{itemize}

\begin{concept}{Saturaci\'on de la Radiancia}
La saturaci\'on ocurre porque al aumentar $\tau$, se cumplen simult\'aneamente:
\begin{itemize}
    \item M\'as part\'iculas dispersan luz $\to$ m\'as fuente de radiancia.
    \item M\'as part\'iculas aten\'uan la luz $\to$ menos fotones llegan a cada
          part\'icula y menos escapan al TOA.
\end{itemize}
El resultado neto es una funci\'on c\'oncava que eventualmente satura en un
valor l\'imite que depende de $\omega_0$, $g$ y $A_s$.
\end{concept}


% ======================================================================
\section{Acoplamiento entre Par\'ametros Recuperados}
\label{sec:interp_coupling}
% ======================================================================

El estudio de recuperaci\'on de 3 par\'ametros revela que el SSA y $g$ tienen
mayor incertidumbre que el AOD.  Esto se explica por el \textbf{acoplamiento
parcial} entre estos par\'ametros en la se\~nal de radiancia:

\subsection{Degeneraci\'on AOD--SSA}

Considerar dos configuraciones:
\begin{align}
    \text{(A)} &\quad \tau = 0.5,\; \omega_0 = 0.9 \implies \tau_{\text{scat}} = 0.45 \\
    \text{(B)} &\quad \tau = 0.45,\; \omega_0 = 1.0 \implies \tau_{\text{scat}} = 0.45
\end{align}

Ambas producen el mismo espesor \'optico de dispersi\'on $\tau_{\text{scat}} =
\omega_0\,\tau$, por lo que generan radiancias similares en dispersi\'on simple.
La diferencia aparece solo en ordenes superiores de dispersi\'on, donde (A) pierde
fotones por absorci\'on en cada evento.

El Jacobiano captura esto: $\partial I / \partial\omega_0$ y
$\partial I / \partial\tau$ tienen perfiles angulares similares (ambos afectan
la magnitud total de la radiancia), lo que produce una alta correlaci\'on
en la covarianza posterior $\hat{\mathbf{S}}$.

\subsection{Degeneraci\'on SSA--$g$}

Un aumento en $g$ concentra la dispersi\'on en la direcci\'on frontal, reduciendo
la radiancia lateral y de retrodispersi\'on.  Un aumento simult\'aneo en $\omega_0$
incrementa el n\'umero total de eventos de scattering, compensando parcialmente.
El averaging kernel $\mathbf{A}$ t\'ipicamente muestra DFS $\approx 2.5$--$2.8$
para 3 par\'ametros, indicando que la informaci\'on observacional no es suficiente
para resolver completamente los tres par\'ametros de forma independiente.


% ======================================================================
\section{Modelamiento para la Comparaci\'on AERONET}
\label{sec:interp_aeronet}
% ======================================================================

\subsection{Cadena de Simulaci\'on}

La comparaci\'on con datos AERONET sigue la cadena:

\begin{equation*}
    \underbrace{\tau_{440}^{\text{AERONET}},\;
    \omega_{0,440}^{\text{AERONET}},\;
    g_{440}^{\text{AERONET}}}_{\text{Datos observados}}
    \xrightarrow{\text{AtmosphereProfile}}
    \underbrace{(\tau, \omega_0, g, A_s, \theta_0)}_{\text{Perfil de 1 capa}}
    \xrightarrow{\text{DISORT}}
    \underbrace{R = F_\uparrow / F_{\text{in}}}_{\text{Reflectancia simulada}}
\end{equation*}

Este enfoque es \textbf{forward modeling}: no se recuperan par\'ametros, sino que
se usan los par\'ametros medidos por AERONET como entrada al modelo directo para
predecir qu\'e reflectancia TOA observar\'ia un sat\'elite.

\subsection{Hip\'otesis y Limitaciones}

\begin{enumerate}
    \item \textbf{Atm\'osfera de 1 capa}: Se modela toda la columna atmosf\'erica
          como una capa homog\'enea, ignorando el perfil vertical real.  Esto es
          una aproximaci\'on razonable para AOD total pero pierde informaci\'on
          sobre la distribuci\'on vertical del aerosol.

    \item \textbf{Superficie Lambertiana con $A_s = 0.1$}: AERONET no reporta
          albedo de superficie.  El valor $0.1$ es representativo de vegetaci\'on
          en el visible pero no se ajusta por d\'ia o estaci\'on.

    \item \textbf{Datos faltantes}: Solo 16 de los 87 d\'ias de AOD tienen datos
          simultan\'eos de SSA y $g$ de inversi\'on AERONET.  Para los d\'ias sin
          inversi\'on, se usan valores por defecto ($\omega_0 = 0.92$, $g = 0.70$),
          introduciendo incertidumbre adicional.

    \item \textbf{Monocrom\'atico a 440\,nm}: Se usa una sola longitud de onda.
          En realidad, los sat\'elites observan en m\'ultiples bandas.
\end{enumerate}

\subsection{Interpretaci\'on de la Relaci\'on AOD--$R$}

La relaci\'on observada entre AOD y $R$ simulado confirma el comportamiento
f\'isico esperado:

\begin{equation}
    R(\tau) \approx R_0 + \frac{\omega_0\,\bar{P}}{4}\left(1 - e^{-c\,\tau}\right)
    \qquad \text{(aproximaci\'on emp\'irica)}
\end{equation}

donde $R_0 = A_s\,e^{-2\tau/\mu_0}$ es la contribuci\'on de superficie (que decrece
con $\tau$), $\bar{P}$ es un promedio efectivo de la funci\'on de fase, y $c$ es una
constante que depende de la geometr\'ia.  La competencia entre estos t\'erminos
produce:
\begin{itemize}
    \item $R$ creciente con $\tau$ para AOD bajo (domina el scattering).
    \item Curvatura y eventual saturaci\'on para AOD alto.
    \item Dependencia del color de los puntos (SSA) confirma que aerosoles
          m\'as absorbentes producen menor $R$.
\end{itemize}


% ======================================================================
\section{Balance Energ\'etico: Verificaci\'on de Consistencia}
\label{sec:interp_energy}
% ======================================================================

La Ec.~\eqref{eq:energy_balance} proporciona una verificaci\'on fundamental del
solver.  El an\'alisis revela comportamientos f\'isicos coherentes:

\subsection{Dependencia con SSA (a $A_s = 0.1$ fijo)}

\begin{itemize}
    \item $\omega_0 = 0$: $F_{\text{abs,atm}}$ es m\'aximo porque toda interacci\'on
          con aerosol es absorci\'on pura.  $F_\uparrow^{\text{TOA}} \approx 0$ excepto
          por la reflexi\'on de superficie del haz directo residual.

    \item $\omega_0 = 1$: $F_{\text{abs,atm}} = 0$ exactamente.  Toda la energ\'ia se
          redistribuye entre $F_\uparrow^{\text{TOA}}$ y $F_{\text{abs,sfc}}$.

    \item Transici\'on: Al aumentar $\omega_0$ de 0 a 1, $F_{\text{abs,atm}}$ decrece
          monotonamente y $F_\uparrow^{\text{TOA}}$ crece, reflejando la mayor eficiencia
          de scattering hacia el espacio.
\end{itemize}

\subsection{Dependencia con Albedo de Superficie (a $\omega_0 = 0.9$ fijo)}

\begin{itemize}
    \item $A_s = 0$: La superficie absorbe todo el flujo descendente.
          $F_{\text{abs,sfc}}$ es m\'aximo.

    \item $A_s = 1$: La superficie no absorbe nada ($F_{\text{abs,sfc}} = 0$).
          Toda la energ\'ia descendente es reflejada, contribuyendo a
          $F_\uparrow^{\text{TOA}}$.

    \item El flujo que escapa al espacio $F_\uparrow^{\text{TOA}}$ crece
          r\'apidamente con $A_s$ porque la superficie act\'ua como una fuente
          adicional de radiaci\'on ascendente.
\end{itemize}


% ======================================================================
\section{Estudio Param\'etrico: Dependencias F\'isicas Clave}
\label{sec:interp_param_study}
% ======================================================================

Los barridos param\'etricos sistem\'aticos revelan las siguientes dependencias:

\subsection{Barrido de SSA ($\omega_0$)}

La radiancia TOA crece monotonamente con $\omega_0$ porque:
\begin{equation}
    \frac{\partial I^{\text{TOA}}}{\partial\omega_0} > 0
    \qquad \text{siempre}
\end{equation}

F\'isicamente, aumentar $\omega_0$ convierte absorci\'on en dispersi\'on,
incrementando la probabilidad de que los fotones escapen hacia el sat\'elite.
El efecto es m\'as pronunciado a AOD alto (m\'as interacciones por fot\'on).

\subsection{Barrido de AOD ($\tau$)}

La radiancia TOA crece con $\tau$ pero con tasa decreciente (curva c\'oncava):
\begin{equation}
    \frac{\partial I}{\partial\tau} > 0, \qquad
    \frac{\partial^2 I}{\partial\tau^2} < 0
    \qquad \text{(para $\tau$ moderado a alto)}
\end{equation}

La primera derivada positiva indica que m\'as aerosol produce m\'as scattering.
La segunda derivada negativa indica saturaci\'on por atenuaci\'on competitiva.

\subsection{Barrido de $g$ (Asimetr\'ia)}

Aumentar $g$ reduce la radiancia en las direcciones de \emph{retrodispersi\'on}
(backscatter) porque la dispersi\'on se concentra hacia adelante:
\begin{equation}
    P(\cos\Theta) = \frac{1-g^2}{(1+g^2-2g\cos\Theta)^{3/2}}
    \xrightarrow{g \to 1} \delta(\cos\Theta - 1)
\end{equation}

Para un sat\'elite observando en la direcci\'on de retrodispersi\'on
($\Theta \approx 180° - \theta_0$), la radiancia \emph{disminuye} con $g$.
Para \'angulos de observaci\'on cercanos a forward scattering, la radiancia
\emph{aumenta} con $g$.

\subsection{Barrido de $\theta_0$ (\'Angulo Cenital Solar)}

Al aumentar $\theta_0$ (sol m\'as bajo):
\begin{equation}
    F_{\text{in}} = F_0\cos\theta_0 \xrightarrow{\theta_0 \to 90°} 0
\end{equation}

El flujo incidente disminuye, pero simult\'aneamente:
\begin{equation}
    \text{camino \'optico oblicuo} = \frac{\tau}{\cos\theta_0}
    \xrightarrow{\theta_0 \to 90°} \infty
\end{equation}

El primer efecto domina para la radiancia absoluta (que disminuye), pero la
\emph{reflectancia} $R = F_\uparrow/(F_0\mu_0)$ puede aumentar porque los
fotones tienen m\'as oportunidades de ser dispersados antes de escapar.


% ======================================================================
\section{Conclusiones}
% ======================================================================

En este trabajo se implement\'o desde cero un solver de ordenadas discretas
para la ecuaci\'on de transferencia radiativa en atm\'osferas plano-paralelas,
siguiendo la metodolog\'ia DISORT documentada por Stamnes et al.\ (2000).
El modelo directo resuelve la RTE discretizada con cuadratura Double-Gauss,
emplea la formulaci\'on reducida de autovalores $(N/2)\times(N/2)$ de
Stamnes \& Swanson (1981) que garantiza autovalores reales por construcci\'on,
e incorpora la estabilizaci\'on num\'erica de Stamnes \& Conklin (1984) y el
escalamiento delta-M de Wiscombe (1977) para funciones de fase fuertemente
asim\'etricas. Adem\'as, se implement\'o la integraci\'on de la funci\'on fuente
(Ecs.~35a,b del reporte) que permite obtener intensidades en \'angulos
arbitrarios sin estar restringidos a los nodos de cuadratura.

La atm\'osfera se model\'o con la Est\'andar US-1976 de 20 capas, incluyendo
dispersi\'on de Rayleigh y propiedades \'opticas de aerosol obtenidas
mediante teor\'ia de Mie (Bohren \& Huffman, 1983) integrada sobre
distribuciones lognormales de tama\~no. Se definieron cinco presets de
aerosol (continental, urbano, mar\'itimo, polvo des\'ertico y quema de biomasa)
y la distribuci\'on vertical se model\'o como un perfil exponencial con
altura de escala $H = 2$\,km.

Para el problema inverso se implementaron dos m\'etodos de regularizaci\'on:
Tikhonov con Levenberg-Marquardt, que minimiza una funci\'on de costo
regularizada mediante un residuo aumentado con el algoritmo Trust Region
Reflective, y Optimal Estimation (Rodgers, 2000), que formula la inversi\'on
en un marco bayesiano con iteraci\'on de Gauss-Newton y proporciona
diagn\'osticos completos como la covarianza posterior, la averaging kernel
y los grados de libertad por se\~nal (DFS). Ambos m\'etodos fueron validados
en los Niveles~1 y~2 del modelo.

El proyecto se estructur\'o en tres niveles de complejidad creciente.
En el \textbf{Nivel~1} se trabaj\'o con una atm\'osfera de una sola capa
y un \'unico par\'ametro a recuperar (AOD), lo que permiti\'o observar
la relaci\'on no lineal entre el espesor \'optico y la radiancia TOA:
a medida que el AOD aumenta, la reflectancia crece pero con tasa
decreciente, de modo que la sensibilidad del observable disminuye
en atm\'osferas \'opticamente gruesas. Se verific\'o que ambos m\'etodos
de inversi\'on recuperan el AOD verdadero con error inferior a $0.01\%$
en ausencia de ruido, y que con a priori d\'ebil ($\mathbf{S}_a$ grande)
el resultado depende exclusivamente de los datos.

En el \textbf{Nivel~2} se pas\'o a una atm\'osfera realista de 20 capas
con aerosol continental a 550\,nm y 36 streams, donde se recuperaron
simult\'aneamente los tres par\'ametros $\mathbf{x} = [\text{AOD},\,\omega_0,\,g]^T$.
Se observ\'o que el AOD es el par\'ametro mejor recuperado, con
RMSE $= 0.089$ sobre 20 escenarios de validaci\'on, mientras que SSA
y $g$ presentan mayor incertidumbre (RMSE de 0.064 y 0.043
respectivamente). Esto se debe a la degeneraci\'on entre $\omega_0$ y $g$:
ambos controlan cu\'anta luz retorna hacia el sensor, de manera que un
aumento en SSA puede ser parcialmente compensado por un aumento en $g$.
La comparaci\'on entre LM+Tikhonov y Optimal Estimation mostr\'o que ambos
convergen a soluciones pr\'acticamente id\'enticas, pero OE ofrece
la ventaja de proporcionar barras de error formales y la averaging kernel
que cuantifica la resoluci\'on de la inversi\'on.

Los estudios param\'etricos permitieron observar el comportamiento f\'isico
esperado en todos los casos: al aumentar $\omega_0$ la reflectancia crece
porque m\'as fotones son dispersados en lugar de absorbidos; al aumentar
$g$ la reflectancia disminuye porque la dispersi\'on se concentra hacia
adelante y menos luz retorna al espacio; y al aumentar el albedo de
superficie, la contribuci\'on del fondo se suma a la radiancia atmosf\'erica.
Se verific\'o la conservaci\'on de energ\'ia exacta
($F_{\text{in}} = F_\uparrow^{\text{TOA}} + F_{\text{abs,sfc}} + F_{\text{abs,atm}}$)
en 13 escenarios, y se confirm\'o que los l\'imites f\'isicos extremos se
reproducen correctamente: reflectancia $R = 1$ para dispersi\'on
conservativa con superficie perfectamente reflectante, y $R = 0$ para
absorci\'on total con superficie negra.

Sin embargo, se identific\'o una limitaci\'on num\'erica importante: para
par\'ametros de asimetr\'ia altos ($g \geq 0.9$), las radiancias presentan
oscilaciones espurias causadas por la truncaci\'on de la expansi\'on de
Legendre de la funci\'on de fase. A $g = 0.95$, la funci\'on de
Henyey-Greenstein requiere aproximadamente 160 t\'erminos de Legendre para
converger, pero el solver utiliza $l_{\max} = N = 36$ t\'erminos (ligado
al n\'umero de streams). El escalamiento delta-M reduce la fracci\'on truncada
pero el par\'ametro de asimetr\'ia escalado ($g' = 0.94$) sigue siendo alto,
y las oscilaciones de tipo Gibbs persisten. Los flujos integrados no se
ven afectados (la cuadratura promedia las oscilaciones), pero el campo
angular s\'i las exhibe. La correcci\'on est\'andar para este problema
(single-scatter correction de Nakajima \& Tanaka, 1988) no fue implementada.

En el \textbf{Nivel~3} se trabaj\'o con datos reales de la estaci\'on
SURFRAD Bondville, que mide irradiancias de banda ancha en superficie
(GHI, DNI, DHI) con piranómetros. La idea central fue explotar la
fracci\'on difusa $d_f = \text{DHI}/\text{GHI}$ como observable: en
cielo despejado, la fracci\'on difusa aumenta con el AOD porque m\'as
fotones son dispersados por los aerosoles antes de llegar al suelo.
Se construy\'o un modelo de 3 bandas (400, 550 y 800\,nm) ponderado por
el espectro solar para hacer el c\'alculo DISORT monocrom\'atico comparable
con las mediciones de banda ancha, y se aplic\'o escalamiento de \AA ngstr\"om
para trasladar el AOD entre longitudes de onda. La inversi\'on se
resolvi\'o por bisecci\'on, aprovechando que $d_f(\text{AOD})$ es una
funci\'on mon\'otona creciente, y el AOD recuperado se valid\'o contra
mediciones independientes de AERONET (fotometr\'ia solar directa)
colocalizadas en la misma estaci\'on.

Este nivel demostr\'o que es posible recuperar AOD a partir de mediciones
puramente de flujo (sin informaci\'on angular), aunque con limitaciones:
el modelo de 3 bandas es una aproximaci\'on gruesa del espectro solar real,
no se incluye absorci\'on gaseosa, y el filtrado de nubes basado en
$d_f < 0.6$ y DNI $> 100$\,W/m$^2$ puede admitir cirros delgados.

Otras limitaciones generales del modelo incluyen: la geometr\'ia
plano-paralela, que ignora la curvatura terrestre y la heterogeneidad
horizontal; la resoluci\'on \'unicamente del modo azimutal $m = 0$, que
produce radiancias promediadas azimutalmente pero no captura la dependencia
con $\phi$; el uso de la funci\'on de fase de Henyey-Greenstein de un
solo par\'ametro, que no reproduce la estructura fina de las funciones de
fase reales; la superficie Lambertiana, que no modela la reflectancia
bidireccional (BRDF) de superficies reales ni la reflexi\'on especular
sobre oc\'eanos; y la ausencia de emisi\'on t\'ermica y de absorci\'on
gaseosa, que restringe la validez del modelo al espectro solar
($\lambda < 3\;\mu$m) y a ventanas atmosf\'ericas libres de l\'ineas
de absorci\'on.

A pesar de estas limitaciones, el modelo constituye una implementaci\'on
funcional y verificada de la cadena completa de teledecci\'on de aerosoles:
desde la microf\'isica de las part\'iculas (teor\'ia de Mie) hasta la
recuperaci\'on de par\'ametros \'opticos a partir de observaciones, pasando
por la transferencia radiativa multicapa con todos los elementos
num\'ericos necesarios para obtener resultados f\'isicamente consistentes.


% ======================================================================
% REFERENCES
% ======================================================================
\begin{thebibliography}{99}

\bibitem{stamnes1988}
Stamnes, K., Tsay, S.-C., Wiscombe, W., \& Jayaweera, K. (1988).
Numerically stable algorithm for discrete-ordinate-method radiative transfer
in multiple scattering and emitting layered media.
\emph{Applied Optics}, 27(12), 2502--2509.

\bibitem{disortreport}
Stamnes, K., Tsay, S.-C., Wiscombe, W., \& Laszlo, I. (2000).
DISORT, a General-Purpose Fortran Program for Discrete-Ordinate-Method
Radiative Transfer in Scattering and Emitting Layered Media:
Documentation of Methodology (v1.1).

\bibitem{rodgers2000}
Rodgers, C. D. (2000).
\emph{Inverse Methods for Atmospheric Sounding: Theory and Practice}.
World Scientific.

\bibitem{chandrasekhar1960}
Chandrasekhar, S. (1960).
\emph{Radiative Transfer}.
Dover Publications.

\bibitem{bohren1983}
Bohren, C. F., \& Huffman, D. R. (1983).
\emph{Absorption and Scattering of Light by Small Particles}.
Wiley.

\bibitem{bodhaine1999}
Bodhaine, B. A., Wood, N. B., Dutton, E. G., \& Slusser, J. R. (1999).
On Rayleigh optical depth calculations.
\emph{J. Atmos. Ocean. Tech.}, 16, 1854--1861.

\bibitem{wiscombe1977}
Wiscombe, W. J. (1977).
The delta-M method: Rapid yet accurate radiative flux calculations for
strongly asymmetric phase functions.
\emph{J. Atmos. Sci.}, 34, 1408--1422.

\bibitem{thomas1999}
Thomas, G. E., \& Stamnes, K. (1999).
\emph{Radiative Transfer in the Atmosphere and Ocean}.
Cambridge University Press.

\bibitem{stamnes1984}
Stamnes, K., \& Conklin, P. (1984).
A new multi-layer discrete ordinate approach to radiative transfer in
vertically inhomogeneous atmospheres.
\emph{J. Quant. Spectrosc. Radiat. Transfer}, 31, 273--282.

\bibitem{holben1998}
Holben, B. N., Eck, T. F., Slutsker, I., et al. (1998).
AERONET --- A federated instrument network and data archive for
aerosol characterization.
\emph{Remote Sens. Environ.}, 66, 1--16.

\bibitem{augustine2000}
Augustine, J. A., DeLuisi, J. J., \& Long, C. N. (2000).
SURFRAD --- A national surface radiation budget network for
atmospheric research.
\emph{Bull. Amer. Meteor. Soc.}, 81, 2341--2358.

\bibitem{stamnes1982}
Stamnes, K. (1982).
On the computation of angular distributions of radiation in
planetary atmospheres.
\emph{J. Quant. Spectrosc. Radiat. Transfer}, 28, 47--51.

\bibitem{sykes1951}
Sykes, J. B. (1951).
Approximate integration of the equation of transfer.
\emph{Mon. Not. R. Astron. Soc.}, 111, 377--386.

\bibitem{stamnes1981}
Stamnes, K., \& Swanson, R. A. (1981).
A new look at the discrete ordinate method for radiative transfer
calculations in anisotropically scattering atmospheres.
\emph{J. Atmos. Sci.}, 38, 387--399.

\end{thebibliography}

\end{document}
